% Français 12357 — lecture modernisée du volume entier.
%
% Pass: Opus 5.5 (claude-opus-5-5), 2026-09-22 — first pass, unchecked against
% the views by a human, and an interpretation rather than a record. Where this
% file and the transcriptions differ in substance, the transcriptions
% (batch-01.fr.tex, batch-02.fr.tex, batch-03.fr.tex) are the record.
%
% Sources read: the three batch transcriptions of this volume, in full, twice —
% once for the mathematics, once for the apparatus — and nothing else. No image
% was consulted. The Harmonie universelle, the book every remark is keyed to,
% was deliberately left shut, and so were Galileo's dialogues, Mersenne's
% other books, Fabre, Faber, Le Fèvre, Guldin and every modern edition: every
% page reference to the Harmonie is given as the copy gives it and none has
% been checked against the book. Where the transcription has a gap, this
% reading says so. All three batches were transcribed on Opus 5.5; together
% they cover views 1–56, the whole volume.
%
% What the volume is. A COPY, in one anonymous seventeenth-century hand, of
% the remarks Mersenne had written in the margins and on the blank leaves of
% his own Harmonie universelle — so says the copy's title (v. 8), and nothing
% on the leaves contradicts it except that a few remarks name « le P.
% Mersenne » in the third person (v. 10, 11, 45), which this reading records
% without drawing a conclusion. Bound in front of the copy: three leaves on
% music in another, older-looking hand, with no reference to the book (v. 3–7).
% Bound at the end: a separate Latin piece, the Paris academy's objections to
% Galileo's dialogues addressed to his friends, dated 1 July 1643 (v. 53–55).
%
% The plan, which follows the treatises the copy's running titles name:
%   v. 1–56   What the volume contains
%   v. 3–7    The musical leaves (plate of 27 strings; monochord; 53 commas)
%   v. 8–9    Blank leaves before the book: circle, square, Fabre, a square root
%   v. 9–11   « Seconde observation »; Book I « de la nature des sons »
%   v. 11–16  Book II « du mouvement des corps »: the spiral of fall and its
%             8/15; not a circle; helices of any degree; odd numbers; media;
%             compound pendulums
%   v. 16–19  Book III « du mouvement »: strings, weights, breaking, missiles
%   v. 19     Book I « de la voix »
%   v. 19–24  Book II « des chants »: permutations and combinations
%   v. 24–28  Books « des consonances »: superparticular divisions, tables
%   v. 28–37  Books « des instruments »: strings, quarter-tones, flutes, bells
%   v. 37–40  « De l'utilité de l'harmonie »: air, mirrors, rational ellipses
%   v. 40–44  « Nouvelles observations »: jets, Beaugrand, temperament,
%             cycloid, amicable and multiply perfect numbers
%   v. 45–52  « Au papier blanc après tous les livres »: magic squares,
%             divisors, primality, sums, dice
%   v. 53–55  The Paris academy to Galileo's friends, 1643
%   v. 1–56   What the leaves do not say
%
% Notation translated, each with a footnote at its first use: « raison
% doublée / sous-doublée » (square / square root of a ratio); « sesquialtère »,
% « sesquioctave », « surparticulière » ((n+1)/n); « partie aliquote » (proper
% divisor); « puissance triangulaire » (figurate number, a binomial
% coefficient); « l'ordre de n », « la combinaison de l'ordre » (n!);
% « quarré quarré », « Q. C. » (fourth, fifth power); « hexagone central »
% (centred hexagonal number); « enceinte » (border of a bordered magic
% square); « filet isochrone » (length of the equivalent simple pendulum);
% « roulette, trochoïde » (cycloid); « hélice » (spiral); « graue » (heavy
% body); « seconde minute » (second of time); the double stroke « ∥ » of v. 13
% (equality, or « est à … comme »); the units toise = 6 pieds, pied = 12
% pouces, pouce = 12 lignes; livre = 16 onces, once = 8 gros, gros = 72 grains.
%
% Every computation the leaves give was redone, and the recomputations are this
% reading's own. Where the leaf is loose, elliptical or wrong this reading
% states what is true and footnotes what the leaf has. Modern names are
% attached where the statement coincides, and footnoted as later than the leaf.
%
% Folios. No \folio{} in any transcription: an ink series 3–26 runs on the
% rectos and matches the librarian's « 26 feuillets », but it is not the pale
% pencil this project takes for the library's foliation, and no human has
% looked. Nothing here is cited by folio; views only, by \pagerange{}{}.
%
% Dates. Only those written on the leaves: « le dernier jour de may 1636 »
% (v. 39, the year doubtful), « l'année 1638 » (v. 19), « le 4e Juin 1638 »
% (v. 41), a letter of Fermat of April 1640 (v. 46, the day doubtful), and
% 1 July 1643 (v. 55). \dating{} is the catalogue's. No date is inferred for the
% copy; the one inference ventured on the remarks (that some cite books
% printed after 1636) is marked as such where it is made.
\documentclass[11pt,a4paper]{article}
% The preamble shared by every transcription of the mathematicians' manuscripts in Gallica.
%
% It rests on one idea: **uncertainty compounds, it does not resolve.** A
% transcription that smooths over a doubtful word has destroyed the only thing
% it was made for — a reader must be able to tell, at every point, what was on
% the page and what was supplied. The apparatus macros below are therefore the
% critical apparatus, not decoration.
%
% The same commands are recognised by `scripts/render.mjs`, which produces the
% site's reading view, and by `scripts/tei.mjs`. A macro added here must be
% added there too, or rendering fails loudly — which is the wanted behaviour.
%
% Comments here are English, like the rest of the repository. The documents
% this preamble sets are French, and that is deliberate: see the skills.

\ProvidesPackage{germain}[2026/09/15 Transcriptions of mathematicians' manuscripts in Gallica]

\usepackage{amsmath,amssymb,amsthm}
\usepackage{mathrsfs}
% Kept from the parent preamble so the renderer's diagram support stays
% available; these manuscripts draw no commutative diagrams, and nothing here uses it.
\usepackage{tikz-cd}
\usepackage[english,french]{babel}
\usepackage{fontspec}

% --- Greek ------------------------------------------------------------------
% The text face stays fontspec's default, Latin Modern, which has no Greek:
% XeTeX drops every glyph it lacks with a line in the log and nothing on the
% page, and the Greek words the manuscripts quote vanished from the PDFs. So
% the Greek and Greek Extended blocks get a XeTeX character class of their own,
% and entering or leaving a run of it switches the family to CMU Serif — the
% same Computer Modern design, with polytonic Greek — and back. Only the
% family changes: a Greek word in \textit or \textbf keeps its shape and series.
% The transitions to and from every other class carry the tokens that class
% has towards an ordinary letter, so French spacing before « ; » or inside
% « » still holds after a Greek word. No grouping, so no brace or \endgroup
% can come out unbalanced; maths is left alone, since XeTeX inserts nothing
% there. Kept identical in `exercices.sty`.
\makeatletter
\newfontfamily\germainGreekfont{cmun}[
  Extension=.otf, NFSSFamily=germaingreek,
  UprightFont=*rm, ItalicFont=*ti, SlantedFont=*sl,
  BoldFont=*bx, BoldItalicFont=*bi, BoldSlantedFont=*bl]
\newXeTeXintercharclass\germain@greek
\def\germain@greekrange#1#2{%
  \@tempcnta=#1\relax
  \loop
    \XeTeXcharclass\@tempcnta=\germain@greek
    \ifnum\@tempcnta<#2\advance\@tempcnta\@ne\repeat}
\germain@greekrange{"0370}{"03FF}
\germain@greekrange{"1F00}{"1FFF}
\def\germain@greekprev{\rmdefault}
\def\germain@greekin{%
  \edef\germain@greekprev{\f@family}\fontfamily{germaingreek}\selectfont}
\def\germain@greekout{\fontfamily{\germain@greekprev}\selectfont}
\def\germain@greekjoin{%
  \XeTeXinterchartokenstate=\@ne
  \@tempcnta=\z@
  \loop
    \ifnum\@tempcnta=\germain@greek\else
      \XeTeXinterchartoks\germain@greek\@tempcnta=\expandafter{%
        \expandafter\germain@greekout\the\XeTeXinterchartoks\z@\@tempcnta}%
      \XeTeXinterchartoks\@tempcnta\germain@greek=\expandafter{%
        \the\XeTeXinterchartoks\@tempcnta\z@\germain@greekin}%
    \fi
    \ifnum\@tempcnta<4095\advance\@tempcnta\@ne\repeat}
% babel-french sets its own classes, and saves and restores the interchar
% state, each time a language is selected: join again after it does.
\addto\extrasfrench{\germain@greekjoin}
\addto\extrasenglish{\germain@greekjoin}
\AtBeginDocument{\germain@greekjoin}
\makeatother

\usepackage{geometry}
\geometry{a4paper,margin=25mm,marginparwidth=28mm}
\usepackage{marginnote}
\usepackage{xcolor}
\usepackage[normalem]{ulem}
\setlength{\emergencystretch}{3em}
\usepackage{hyperref}
\hypersetup{colorlinks=true,linkcolor=black,urlcolor=black,pdfborder={0 0 0}}

% --- Batch metadata ---------------------------------------------------------
% Taken as they stand by the rendering script, which shows them at the head of
% the reading view. They fix what the file claims to cover. The macro names are
% the parent project's, so that the scripts stay shared; \folder here means the
% volume's slug — `fr-9115` — and \pages counts Gallica views.

\newcommand{\folder}[1]{\gdef\grothFolder{#1}}
\newcommand{\batch}[1]{\gdef\grothBatch{#1}}
\newcommand{\pages}[2]{\gdef\grothFirst{#1}\gdef\grothLast{#2}}
\newcommand{\foldertitle}[1]{\gdef\grothFoldertitle{#1}}

% The holder's own shelfmark, printed as they write it: « Français 9115 ».
\newcommand{\shelfmark}[1]{\gdef\grothShelfmark{#1}}

% The dating, copied verbatim from the holder's catalogue — never inferred
% here. « XIXe siècle » is the BnF's; a pass that dates a leaf from its content
% says so in a \note{}, as its own inference.
\newcommand{\dating}[1]{\gdef\grothDating{#1}}

% The status notice, printed in the title block and repeated as a diagonal
% watermark on every page of the PDF. The manuscripts are in the public
% domain; what this line declares is not a right but a quality — an
% unverified first pass by a machine. The skills write the line; a file
% without it gets no watermark, so leaving it out is visible in the output.
\newcommand{\watermark}[1]{\gdef\grothWatermark{#1}}

\gdef\grothFolder{?}\gdef\grothBatch{}\gdef\grothFirst{?}\gdef\grothLast{?}%
\gdef\grothFoldertitle{}\gdef\grothShelfmark{}\gdef\grothDating{}\gdef\grothWatermark{}

% --- The résumé -------------------------------------------------------------
\newenvironment{resume}
  {\small\setlength{\parskip}{0.45em}}
  {\par\medskip\hrule\medskip}

% --- Keywords ---------------------------------------------------------------
% In English, closing the modernised reading's résumé: the modern vocabulary
% under which to search for what the leaves construct. The manifest extracts
% them to tag the volume — this line is the single source of the tags.
\newcommand{\keywords}[1]{\par\smallskip\noindent
  {\footnotesize\textsc{Keywords} --- \textit{#1}}\par}

% --- The critical apparatus -------------------------------------------------

% The Gallica view number, in the margin. This is the anchor that lets a
% disagreement over a formula be settled by looking at *one* image rather than
% a volume of 750. The site uses it to turn the facsimile as the transcription
% is scrolled. A view is one scanned image: usually one side of a leaf.
\newcommand{\page}[1]{%
  \marginnote{\footnotesize\color{gray!70}\textsf{v.\,#1}}%
  \label{page:#1}\ignorespaces}

% The BnF's pencilled foliation, where the leaf carries it and a pass can read
% it: « 198r ». This is what the literature cites — « ff. 198r–208v » — and
% the one line that turns a citation into a view. Recorded, never inferred:
% a folio a pass cannot read is not written.
\newcommand{\folio}[1]{%
  \marginnote{\footnotesize\color{gray!50}\textsf{f.\,#1}}\ignorespaces}

% The modernised reading's coarser counterpart to \page{}: which views a
% section is drawn from.
\newcommand{\pagerange}[2]{%
  \marginnote{\footnotesize\color{gray!70}\textsf{v.\,#1--#2}}%
  \ignorespaces}

% Illegible: never guessed.
\newcommand{\ill}{\textcolor{red!60!black}{[\,\ldots\,]}}

% A doubtful reading: the word is given, and flagged as uncertain.
\newcommand{\uncertain}[1]{\uline{#1}}

% An editorial addition — a missing letter, an implied word.
\newcommand{\add}[1]{\textcolor{blue!50!black}{[#1]}}

% Struck out by the writer. What was crossed out is part of the document.
\newcommand{\struck}[1]{\sout{#1}}

% The transcriber's note, on the state of the page or the mathematics.
\newcommand{\note}[1]{\par\smallskip{\small\color{gray!60!black}\textit{[#1]}}\par\smallskip}

% A marginal note of hers, distinct from the body of the text.
\newcommand{\marginal}[1]{\par\smallskip\noindent\hspace{1em}%
  {\small\color{gray!60!black}#1}\par\smallskip}

% --- Title ------------------------------------------------------------------
% The title block is French, like the documents themselves.

\AddToHook{shipout/background}{%
  \ifx\grothWatermark\empty\else
    \begin{tikzpicture}[remember picture, overlay]
      \node[rotate=52, scale=3.4, align=center, text=gray!18,
            font=\sffamily\bfseries]
        at (current page.center) {\grothWatermark};
    \end{tikzpicture}%
  \fi}

\AtBeginDocument{%
  \begin{center}
    {\large\bfseries Manuscrits de mathématiciens (Gallica) ---
      \ifx\grothShelfmark\empty\grothFolder\else\grothShelfmark\fi}\\[2pt]
    {\itshape\grothFoldertitle}\\[4pt]
    {\small vues \grothFirst--\grothLast
      \ifx\grothBatch\empty\else\ (lot \grothBatch)\fi}\\[2pt]
    \ifx\grothDating\empty\else
      {\small\color{gray!60!black}Datation du catalogue : \grothDating}\\[2pt]
    \fi
    \ifx\grothWatermark\empty\else
      {\small\bfseries\color{red!45!gray}\grothWatermark}\\[2pt]
    \fi
    {\footnotesize\color{gray!60!black}Transcription établie d'après les images IIIF de Gallica.
     Source gallica.bnf.fr / Bibliothèque nationale de France.\\
     Les lectures incertaines sont soulignées, les illisibles notées [\,\ldots\,],
     les ajouts entre crochets ; \textsf{v.} numérote les vues de Gallica,
     \textsf{f.} la foliotation au crayon de la BnF.}
  \end{center}
  \vspace{1em}}

\endinput


\folder{fr-12357}
\shelfmark{Français 12357}
\pages{1}{56}
\dating{1601-1700}
\watermark{Lecture automatique\\interprétation non vérifiée}
\foldertitle{« Remarques tirées du livre de l'Harmonie universelle du P. MERSENNE, ainsy qu'il les avoit escrites, de sa main, à la marge et aux feuillets blancs devant et derrière dudit livre ». Français 12357 — lecture modernisée du volume entier}

\begin{document}

\section*{Résumé}

\begin{resume}
Ce volume est la copie, par une main anonyme du XVII\textsuperscript{e}
siècle, des notes que le P. Marin Mersenne avait écrites dans les marges et sur
les feuillets blancs de son propre exemplaire de l'\emph{Harmonie universelle},
le grand traité de musique, d'acoustique et de mécanique qu'il avait publié. La
copie le dit dans son titre : ce sont les remarques « ainsi qu'il les avait
écrites de sa main ». Chaque remarque commence par le renvoi à la page du livre
qu'elle commente — « p. 94 au commencement », « p. 189, prop. XV » — et la copie
suit le livre traité par traité : la nature des sons, le mouvement des corps,
la voix, les chants, les consonances, les instruments, l'utilité de
l'harmonie, puis les « nouvelles observations », enfin ce qu'il avait écrit
« au papier blanc après tous les livres ». Trois feuillets de musique d'une
autre main sont reliés devant, et une pièce latine datée du 1\textsuperscript{er}
juillet 1643 est reliée derrière : les objections que l'« Académie parisienne »
avait écrites contre les dialogues de Galilée et qu'elle adresse, Galilée étant
mort, à ses amis d'Italie.

Le problème que se pose l'auteur des remarques est celui d'un savant qui relit
son livre et le corrige. Il le dit sans détour : « il faut corriger tout ceci »,
« ce nombre me semble faux », « il la faut ôter d'ici comme inutile », « voyez
la 9\textsuperscript{e} proposition de l'Utilité, et corrigez tout ceci ». Ses
corrections portent sur la vitesse du son, le poids de l'air, la chute des
graves, la portée des projectiles, et elles renvoient sans cesse à des
expériences refaites, à des lettres de Descartes, de Fermat, de Galilée. Mais
le volume a aussi ses pages de pure mathématique, et elles sont de premier
ordre. Un corps qui tomberait vers le centre d'une Terre en rotation ne décrit
pas le demi-cercle qu'avait proposé Galilée, mais une spirale ; l'aire de cette
spirale est exactement les $8/15$ du secteur qu'elle parcourt, démontrés à la
manière d'Archimède par des secteurs inscrits et circonscrits et un lemme sur
les sommes de carrés que la marge vérifie en chiffres ; une autre démonstration,
en algèbre, établit que la courbe ne peut pas être un cercle ; et une petite
table sans titre donne, pour les spirales de tout degré, les rapports
$1/3, 8/15, 9/14, 32/45, \dots$ — qui sont exactement $2n^2/((n+1)(2n+1))$.
Ailleurs, on dénombre les mains de douze cartes au piquet, les chants de huit
notes, les mots qu'on peut faire de 23 lettres ; on construit des carrés
magiques de 4 à 22 cases de côté, emboîtés les uns dans les autres ; on
compte les diviseurs d'un nombre, on cherche le moindre nombre qui en a un
nombre donné, on énumère les nombres amiables et les nombres multiparfaits ; on
demande si 1073602561 et 2027651281 sont premiers ; on compte les coups de
deux et de trois dés.

L'idée qui fait tenir ensemble ces pages disparates est celle de l'auteur de
l'\emph{Harmonie} lui-même : tout se mesure et tout se compte, la consonance par
la coïncidence des vibrations, la hauteur d'une corde par le poids qui la
tend, le nombre des chants possibles par la combinatoire. Cette lecture a
refait tous les calculs. La plupart sont justes, et certains le sont à un
degré remarquable : la table des parties aliquotes, les trois paires de
nombres amiables, les six nombres dont la somme des diviseurs est triple du
nombre, la vérification chiffrée du lemme des $8/15$, les angles d'une cloche.
D'autres sont faux, et on le dit à chaque fois : 6679 n'est pas le produit de 23
et 29, le moindre nombre qui a onze parties aliquotes est 60 et non 72, 15000 en
a 39 et non 79, la quarte ne se divise pas en 27 mais en 26 manières en trois
raisons superparticulières, deux dés ne donnent pas 62 coups mais 36. Plusieurs
fautes sont manifestement de copie, comme le 4320 qui tient la place de 40320.

Où cela mène, et quels noms modernes chercher : la chute des corps dans une
Terre en rotation et la mécanique en référentiel tournant ; la quadrature par
exhaustion et les sommes de puissances ; le pendule composé et la longueur
du pendule simple équivalent ; les lois de Mersenne pour les cordes
vibrantes ; les permutations, les combinaisons avec et sans répétition et les
coefficients binomiaux ; les divisions superparticulières, l'intonation juste,
le tempérament égal et la division de l'octave en 53 ; les carrés magiques à
enceintes ; la fonction nombre de diviseurs $\tau$ et la fonction somme des
diviseurs $\sigma$ ; les nombres amiables et la règle de Thābit ibn Qurra ; les
nombres multiparfaits ; la méthode de factorisation de Fermat et les nombres de
Mersenne ; la loi de Torricelli ; la cycloïde ; le calcul des chances aux dés.

\keywords{Mersenne's Harmonie universelle, marginalia, fall towards the centre of a rotating Earth, area of a spiral, quadrature by exhaustion, sums of powers, compound pendulum, equivalent simple pendulum, Mersenne's laws of vibrating strings, tensile strength, permutations, combinations with repetition, binomial coefficients, figurate numbers, inclusion-exclusion principle, superparticular ratios, just intonation, Pythagorean tuning, equal temperament, 53-tone equal temperament, quarter-tone scale, bordered magic squares, divisor function, sum-of-divisors function, amicable numbers, Thabit ibn Qurra's rule, multiply perfect numbers, Fermat's factorization method, Mersenne primes, Torricelli's law, cycloid, dice probabilities, Galileo's Two New Sciences}
\end{resume}

\section*{Ce que ce volume contient}

\pagerange{1}{56}
Cinquante-six vues, une page par vue. La vue 1 est le plat marbré de la
reliure, la vue 2 la garde, qui ne porte que l'ancienne cote « Suppl.\textsuperscript{t}
fr. 902 » et la mention d'un bibliothécaire, « Volume de 26 Feuillets 3 Août
1894 » ; la vue 56 est le verso blanc du dernier feuillet, et la vue 52 ne
montre que la page précédente par transparence. Tout le reste se partage en
trois ensembles.

\begin{itemize}
  \item \textbf{Vues 3 à 7 — trois feuillets de musique} d'une main plus ancienne
    d'aspect que celle de la copie, aux s et aux e bouclés : une grande planche
    dépliante (vue 3), un monocorde divisé (vue 4), un clavier gradué en commas
    (vue 5), deux pages de musique notée à trois et quatre parties (vues 6
    et 7). Aucun renvoi à une page de l'\emph{Harmonie} n'y figure.
  \item \textbf{Vues 8 à 52 — la copie des remarques}, d'une seule main du
    XVII\textsuperscript{e} siècle, sous le titre « Remarques tirées du livre de
    l'harmonie universelle du P. Mersenne ainsi qu'il les avait écrites de sa
    main à la marge et aux feuillets blancs devant et derrière dudit livre ».
  \item \textbf{Vues 53 à 55 — une pièce latine}, de la même main mais d'une
    écriture plus posée : « Academia Parisiensis viros Clarissimos Galilaei
    familiares et amicos lyncaeos precatur\dots », quinze objections numérotées
    aux dialogues de Galilée, datées de Paris le 1\textsuperscript{er} juillet
    1643, suivies d'un tableau des sept modes grecs d'après Doni.
\end{itemize}

La copie suit le livre. Chaque remarque s'ouvre sur son renvoi — la page, puis
le lieu de la page : « au droit des lignes 23 jusqu'à 32 », « lig. 22 après ces
mots (\dots) », « au droit de la figure », « sur le corollaire premier » — et la
pagination repart à chaque traité du livre, que la copie nomme dans sa marge
gauche sous un trait. Les sections de cette lecture suivent ces titres
courants, dans l'ordre de la copie.\footnote{Les titres courants relevés par les
transcriptions sont : « au livre 1\textsuperscript{er} de la nature des sons »
(v. 9), « au livre second du mouvement des corps » (v. 11), « au
3\textsuperscript{me} livre du mouvement » (v. 16), « Du 1\textsuperscript{er}
livre de la voix » (v. 19), « Du livre second des chants » (v. 19),
« \dots 1\textsuperscript{er} des \dots sonances » (v. 24, en partie dans le pli),
« Du livre 3\textsuperscript{me} des consonnances » (v. 25), « Des Instruments
\dots chorde » (v. 28), « \dots ivre second des Instruments » (v. 29),
« \dots ivre 3\textsuperscript{me} des \dots truments a chorde » (v. 30),
« \dots vij\textsuperscript{me} livre instruments de \dots ussion » (v. 33, le
numéro douteux), « De l'utilité de l'harmonie » (v. 37), « \dots nouvelles
\dots servations \dots siques et mathématiques » (v. 40), et, en tête de page,
« Au papier blanc après tous les livres de l'harmonie » (v. 45). Entre deux
titres, l'appartenance d'une page à un traité est celle que donne la copie ;
aucun renvoi n'a été confronté au livre.} Les renvois aux pages de
l'\emph{Harmonie} sont donnés ici tels que la copie les écrit ; cette lecture
n'a pas ouvert le livre et ne les a pas vérifiés.

La copie garde la première personne de l'auteur des remarques : « j'ai donné ma
méthode », « notre version », « j'ai vu graver à Mons », « j'ai trouvé par la
pompe éolipile ». Trois passages pourtant nomment Mersenne à la troisième
personne : un écho « répond une octave plus haut que la voix du P. Mersenne »
(v. 10) ; deux colonnes d'une table donnent les chutes « suivant les
expériences du père Mersenne » (v. 11) ; et, en marge du grand carré magique
(v. 45), la copie corrige trois sommes « selon le père Mersenne, mais il est
faux ». Cette lecture relève ces passages sans en tirer de conclusion sur la
personne qui les a écrits ; elle attribue les remarques à Mersenne parce que
le titre de la copie le fait, et elle n'attribue la copie à personne.

Pourquoi les feuillets ne sont cités que par les vues de Gallica. Une série de
petits chiffres fins, à l'encre, court au coin supérieur droit des rectos, de
« 3 » (vue 6) à « 26 » (vue 55), un par feuillet, et finit sur le compte de
vingt-six feuillets du bibliothécaire ; elle se comporte donc comme une
foliotation. Mais elle est à l'encre et non au crayon, et aucune transcription
ne l'a écrite comme folio en attendant qu'une personne ait regardé les
feuillets. Cette lecture fait de même.\footnote{Autres chiffres relevés par les
transcriptions : un « 331 » au-dessus du titre (v. 8) ; un « (4 » à la fin de la
troisième ligne du titre, là où l'on attendrait le chiffre de la série, mais
d'un trait plus gras (v. 8) ; un « 2 » et un « 3 » isolés au pied des vues 24 et
42, peut-être des signatures de cahier ; un « 1 » ou « i » en tête de la vue 54.
Tous les « p. », « pag. », « pa. » du texte sont des pages du livre imprimé
(ou, aux vues 8 et 9, du livre de Fabre), jamais des pages du manuscrit.}

\section*{Les feuillets de musique reliés devant la copie}

\pagerange{3}{7}
\subsection*{La planche des vingt-sept cordes}

La planche de la vue 3 donne vingt-sept longueurs de cordes en escalier, de
57600 à 28800, c'est-à-dire sur une octave. Rapportées à la plus longue, toutes
sauf deux sont des rapports de l'intonation juste, produits de puissances de
2, 3 et 5 : $55296 = 57600 \cdot \frac{24}{25}$ (le dièse chromatique),
$54000 = 57600\cdot\frac{15}{16}$, $51840 = 57600\cdot\frac{9}{10}$,
$51200 = 57600\cdot\frac{8}{9}$, $48000 = 57600\cdot\frac56$,
$46080 = 57600\cdot\frac45$, $38400 = 57600\cdot\frac23$,
$34560 = 57600\cdot\frac35$, $30720 = 57600\cdot\frac{8}{15}$, et ainsi des
autres. Les deux exceptions sont les douzième et treizième cordes, lues 43500
et 43472 : aucune n'est un rapport simple, et la quarte juste,
$57600 \cdot \frac34 = 43200$, manque à l'escalier. Il est probable que l'une
des deux devait être 43200 ; cette lecture ne devine pas l'autre.\footnote{La
transcription précise que les deux nombres ont été lus « tels qu'ils sont
tracés et vérifiés à fort grossissement ». Mesurées en commas de l'octave
divisée en 53, les vingt-cinq autres cordes tombent toutes à moins de
$0{,}2$ comma d'un entier ; celles-ci tombent à $21{,}47$ et $21{,}52$, entre
la quarte (22) et rien.}

Les deux tables de la planche nomment les intervalles et leurs rapports. Celle
des consonances est juste : tierces mineure $\frac56$ et majeure $\frac45$,
quarte $\frac34$, quinte $\frac23$, sixtes $\frac58$ et $\frac35$, octave
$\frac12$ — le rapport étant celui de la corde aiguë à la grave. Celle des
dissonances l'est aussi pour l'essentiel — comma $\frac{80}{81}$, dièse
$\frac{125}{128}$, demi-tons $\frac{24}{25}$, $\frac{128}{135}$,
$\frac{15}{16}$, tons $\frac{9}{10}$ et $\frac89$, triton $\frac{32}{45}$,
fausse quinte $\frac{45}{64}$ ; la « fausse quarte » $\frac{75}{96}$ est
$\frac{25}{32}$, la « quinte superflue » $\frac{48}{75}$ est $\frac{16}{25}$ —
à deux articles près. La « septième mineure » $\frac{148}{225}$ n'est le rapport
d'aucune septième : son inverse, $1{,}52$, est à peine plus grand qu'une quinte.
Et la « septième majeure » $\frac{10}{18} = \frac59$ est en réalité la septième
mineure $\frac95$. L'escalier lui-même contient les bonnes valeurs : $32400$ et
$32000$ (les septièmes mineures $\frac{9}{16}$ et $\frac59$) et $30720$ (la
septième majeure $\frac{8}{15}$).\footnote{La transcription lit « 148 » en
précisant que « le premier chiffre après 1 est un 4 net ». Le « comma mineur »
$\frac{2025}{2048}$ est ce qu'on appelle aujourd'hui le diaschisma. Trois
claviers sont dessinés sous la planche, « Octave de vingt-sept, de
vingt-cinq, de dix-neuf marches » ; la grille des consonances, qui n'est pas
transcrite case par case, n'est pas lue ici.}

\subsection*{Le monocorde et ses feintes}

La vue 4 divise un monocorde et explique comment placer les « feintes », les
touches noires, à partir des consonances. « Soustraire » et « ajouter » des
raisons, sur ce feuillet, c'est les diviser et les multiplier : ainsi
$A = \frac{2}{3} : \frac{3}{4} = \frac89$, $F = \frac34\cdot\frac34 =
\frac{9}{16}$, do dièse $= \frac45\cdot\frac89 = \frac{32}{45}$, sol dièse
$= \frac{32}{45} : \frac34 = \frac{128}{135}$, fa dièse $= \frac{32}{45}\cdot
\frac34 = \frac{8}{15}$. Les cinq calculs sont justes.\footnote{« Pour
trouver A faut soubtraire la raison harmonique $\frac34$ de D, qui fait
$\frac23$, reste $\frac89$ pour A. » Le signe de la feinte est, sur le
feuillet, la lettre suivie d'un z à queue ; la transcription le rend par un
dièse. Les valeurs $\frac{96}{135}$, $\frac{64}{135}$ et $\frac{16}{45}$ de la
figure sont $\frac{32}{45}$ et les octaves aiguës de $\frac{128}{135}$ et de
$\frac{32}{45}$.}

\subsection*{Le clavier en 53 commas}

La vue 5 attribue à chaque intervalle un nombre entier de « commas », l'octave
en comptant 53 : dièse 3, « lima » 4, demi-ton 5, ton mineur 8, ton majeur 9,
ton « excédant » 10, tierce mineure 14, tierce majeure 17, quarte 22, quinte
31, sixte mineure 36, sixte majeure 39, octave 53. C'est, en langage moderne,
l'approximation des intervalles justes dans la division de l'octave en 53
parties égales : un intervalle de rapport $r$ y vaut $53\log_2 r$ commas, et
l'on trouve $3{,}12$ ; $4{,}07$ ; $4{,}93$ ; $8{,}06$ ; $9{,}01$ ; $9{,}87$ ;
$13{,}94$ ; $17{,}06$ ; $22{,}00$ ; $31{,}00$ ; $35{,}94$ ; $39{,}06$ — tous à
moins d'un septième de comma de l'entier du feuillet. Les compositions que le
feuillet donne sont exactes, en commas comme en rapports : le ton mineur est le
dièse plus le demi-ton ($\frac{25}{24}\cdot\frac{16}{15} = \frac{10}{9}$,
$3 + 5 = 8$), le ton majeur la lima plus le demi-ton ($\frac{135}{128}\cdot
\frac{16}{15} = \frac98$, $4 + 5 = 9$), la sixte mineure deux tons majeurs, un
ton excédant et un ton mineur ($\frac{81}{64}\cdot\frac{256}{225}\cdot
\frac{10}{9} = \frac85$, $36$), l'octave trois tons majeurs, deux mineurs et un
excédant ($\frac{729}{512}\cdot\frac{100}{81}\cdot\frac{256}{225} = 2$,
$53$).\footnote{Le feuillet appelle « ton parfait » le ton $\frac98$, « ton
majeur ou excédant » le rapport $\frac{256}{225}$ et dit du comma qu'il est
« la $\frac{1}{53}$ partie de l'octave » ; le comma syntonique $\frac{81}{80}$
en est en fait $\frac{1}{55{,}8}$. La gamme est dite « inventée par Boetius
pour la réformation de la musique sous Grégoire premier » ; cette lecture
n'en juge pas. Le nom moderne de la division, tempérament à 53 degrés égaux,
est postérieur au feuillet. La table des « concordances et discordances »
marquées $+$ et $-$ n'est pas vérifiée case par case.} Les vues 6 et 7 sont de
la musique notée, non transcrite.

\section*{Les feuillets blancs de devant : le cercle, le carré, et Fabre}

\pagerange{8}{9}
Les premières remarques ne commentent pas l'\emph{Harmonie} : elles étaient,
dit la marge, sur les feuillets blancs placés devant le livre, et elles
discutent la quadrature du cercle d'après un livre de Fabre dont elles citent
les pages 73 à 97.\footnote{Ces pages sont celles de Fabre, non de
l'\emph{Harmonie}, d'après le texte même (« Fabre p. 73 prétend\dots »). Le livre
n'est pas autrement nommé.}

\subsection*{Aires et périmètres, avec $\pi \approx \frac{22}{7}$}

Tous les calculs de cette page prennent $\pi = \frac{22}{7}$. Une sphère de
diamètre 7 a un grand cercle de circonférence 22 et une surface totale 154
(diamètre fois circonférence, soit $\pi d^2$) ; le carré de même périmètre que
le cercle de diamètre 7 a pour côté $5\frac12$. Un cercle de diamètre 14 a pour
aire 154, et le carré de même périmètre, de côté 11, n'a que 121 : les deux aires
sont comme 14 à 11 pour un même tour de 44 pieds. Tout cela est juste.

« Pour faire un rond égal à un carré », on donne 8 au diamètre du cercle et 10
à la diagonale du carré. L'aire du carré est 50 ; celle du cercle, avec
$\frac{22}{7}$, est $16 \cdot \frac{22}{7} = 50\frac27$, et non $50\frac47$
comme l'écrit la copie ; l'excès du cercle est donc $\frac{1}{176}$ de sa
propre aire, et non $\frac{1}{88}$.\footnote{Avec la vraie valeur de $\pi$,
l'aire est $50{,}27$ et l'excès $\frac{1}{189}$ environ. Les deux nombres de la
copie, $50\frac47$ et $\frac{1}{88}$, sont le double de l'excès vrai ; ils
s'accordent entre eux et non avec le calcul.}

« Par le moyen de 70 et 99 multipliés en soi-mêmes l'on a deux nombres en
raison double, à $\frac{1}{4900}$ près. » En effet $99^2 = 9801$ et
$2 \cdot 70^2 = 9800$ : $\left(\frac{99}{70}\right)^2 = 2 + \frac{1}{4900}$.
C'est une solution de l'équation $x^2 - 2y^2 = 1$, qu'on appelle aujourd'hui
équation de Pell-Fermat, et $\frac{99}{70}$ est une réduite du développement de
$\sqrt2$ en fraction continue.\footnote{Les deux noms sont postérieurs au
feuillet. La marge dessine le carré de côté 70 et de diagonale 99.}

\subsection*{Trois valeurs de $\pi$}

La page oppose trois approximations. Celle de Fabre fait égale au quart de la
circonférence circonscrite la droite AG menée d'un sommet A du carré au milieu
G du côté opposé. Si le carré a pour côté $a$, $AG = \frac{\sqrt5}{2}a$, et le
quart de la circonférence circonscrite est $\frac{\pi}{2\sqrt2}a$ : la règle
revient à $\pi = \sqrt{10} \approx 3{,}1623$. La propriété annexe, que GH
(joignant les milieux de deux côtés) égale le rayon AE, est exacte.

Celle du sieur Cornu fait que, le côté du carré inscrit étant 9, le quart de la
circonférence est 10. Le côté inscrit vaut $r\sqrt2$ et le quart de la
circonférence $\frac{\pi r}{2}$ : la règle revient à
$\pi = \frac{20\sqrt2}{9} \approx 3{,}142697$, trop grand de $1{,}1\cdot10^{-4}$.
La remarque ajoute que Cornu fait le quart de la circonférence « trop grand
d'un peu, car le polygone circonscrit de 97 côtés est plus petit que son
cercle supposé ». C'est exact, et de justesse : le périmètre du polygone
régulier de 97 côtés circonscrit au cercle de diamètre 1 vaut
$97\tan\frac{\pi}{97} \approx 3{,}142692$, inférieur de cinq millionièmes à la
valeur de Cornu. Un cercle ne pouvant être plus long que les polygones qui
lui sont circonscrits, la valeur de Cornu est bien trop grande.

La troisième est citée d'un « Denis \dots, conseiller du Roi au bailliage de
Paris » : la circonférence contient le diamètre « trois fois et
$\frac{112}{113}$ d'une septième partie ». Or
$3 + \frac{112}{113}\cdot\frac17 = 3 + \frac{16}{113} = \frac{355}{113}
\approx 3{,}1415929$, la valeur qu'on attribue aujourd'hui à Adriaan Metius,
exacte à $3\cdot10^{-7}$ près.\footnote{Le nom du conseiller est souligné et lu
« auis » ou « anis » par la transcription. L'attribution de $\frac{355}{113}$ à
Metius est un usage moderne ; le feuillet ne la fait pas.}

Enfin, pour l'aire du cercle : « multiplier le diamètre par soi-même, puis par
11, et diviser par 14 » est juste ($\frac{\pi}{4} \approx \frac{11}{14}$) ; mais
« la moitié du diamètre par la circonférence » donne le double de l'aire. Il
faut la moitié du diamètre par la moitié de la circonférence,
$\frac{d}{2}\cdot\frac{C}{2} = \pi r^2$.

\subsection*{Cercles additionnés, anneaux, triangle}

La construction de Guillaume Cambray pour faire un cercle égal à deux cercles
donnés, et la soustraction d'un cercle d'un autre qui laisse un anneau
(« armille ») égal à un cercle, ne sont que le théorème de Pythagore appliqué
aux rayons : $\pi R^2 = \pi r_1^2 + \pi r_2^2$ si et seulement si
$R^2 = r_1^2 + r_2^2$. La remarque le dit elle-même, avec humeur : « ce qu'étant
connu aux enfants par les éléments d'Euclide, je vois que Fabre n'était ni
géomètre ni philosophe. » Les rapports de la page 95 de Fabre sont justes avec
$\frac{22}{7}$ : le cercle inscrit dans un carré est au carré comme
$\frac{\pi}{4} = \frac{11}{14}$, le carré inscrit dans un cercle est au cercle
comme $\frac{2}{\pi} = \frac{7}{11}$.

La réduction d'un carré en triangle équilatéral (p. 93 de Fabre) n'est pas
exacte. Pour un carré de côté $a$, le triangle équilatéral de même aire a pour
demi-côté $a/\sqrt[4]{3} \approx 0{,}7598\,a$. La construction prend pour
demi-côté la distance CE, E étant sur le quart de cercle BEC. Si E est le
milieu de l'arc, $CE = 2a\sin 22{,}5^\circ \approx 0{,}7654\,a$, et le triangle
dépasse le carré d'environ $1{,}5\,\%$ ; si E est, comme le dit à la lettre la
note marginale, à la distance « moitié de BC », c'est-à-dire la
demi-diagonale $a/\sqrt2$, le triangle ne vaut que $\frac{\sqrt3}{2} \approx
0{,}87$ du carré.\footnote{La figure de la transcription ne permet pas de
situer E plus précisément ; les deux lectures sont données.} La règle de la page
97 — l'aire du triangle équilatéral est le carré du côté multiplié par 13 et
divisé par 30, sa hauteur le côté multiplié par 13 et divisé par 15 — revient à
$\sqrt3 \approx \frac{26}{15} = 1{,}7333$ au lieu de $1{,}7321$. Pour le côté 8
elle donne $27\frac{11}{15}$ et $6\frac{14}{15}$, justes dans cette
approximation ; les valeurs vraies sont $27{,}71$ et $6{,}93$.

\subsection*{Le carré de 3,74586}

La deuxième page blanche porte le carré de $3{,}74586$ posé chiffre à chiffre :
une ligne pour les carrés des chiffres ($9$, $49$, $16$, $25$, $64$, $36$), une pour les
doubles produits de chaque chiffre avec le précédent ($42$, $56$, $40$, $80$, $96$), une
avec l'avant-précédent ($24$, $70$, $64$, $60$), et ainsi de suite jusqu'au double
produit du premier et du dernier ($36$). C'est le développement
\[
\Bigl(\sum_i d_i\,10^{-i}\Bigr)^2 = \sum_i d_i^2\,10^{-2i}
  + 2\sum_{i<j} d_i d_j\,10^{-i-j},
\]
rangé par écart $j - i$. Toutes les lignes sont justes et le total,
$14{,}0314671396$, est bien le carré de $3{,}74586$.

\section*{La « seconde observation » et le livre de la nature des sons}

\pagerange{9}{11}
Les remarques sur la « seconde observation, qui est après la table des
matières » reprennent la théorie de la consonance par la coïncidence des
vibrations. Deux cordes à la quarte ($4 : 3$) ne s'unissent qu'une fois sur
quatre battements de l'aiguë, dont trois restent « désunis » ; à la quinte
($3 : 2$), deux sur trois. Si l'on superpose trois cordes formant une octave
divisée par une quarte en bas, leurs fréquences sont comme $3 : 4 : 6$ et
toutes trois ne s'unissent qu'au sixième battement de l'aiguë ; si la quinte
est en bas, elles sont comme $2 : 3 : 4$ et s'unissent au quatrième. La
remarque est exacte.

Au livre de la nature des sons, trois nombres reviennent. Le son parcourt 1380
pieds, soit 230 toises, par seconde, et le coup de canon tiré de l'Arsenal
s'entend aux Minimes de Vincennes en 20 secondes, « qu'on voie s'il y a 4600
toises » : l'arithmétique est juste ($230 \cdot 20 = 4600$), la valeur ne l'est
pas. La vitesse du son dans l'air est d'environ $340$ mètres par seconde, soit
à peu près $1050$ pieds de roi.\footnote{Le pied de roi vaut $32{,}48$ cm, la
toise six pieds. La copie dit que la vitesse est la même « fort ou faible »,
en tout temps et en tout lieu : c'est vrai de l'intensité, à peu près vrai du
lieu, et faux de la température, dont la vitesse dépend.} L'air est « 1700 fois
plus léger que l'eau », contre les 400 de Galilée ; ailleurs le volume dira
1300, 1356, 1380, et Descartes 145. La valeur vraie, au niveau de la mer et à
température ordinaire, est d'un peu plus de 800. Enfin, « les cercles de l'air
sont 1380 fois plus vites que ceux de l'eau » (p. 67) confond le rapport des
densités avec celui des vitesses : le son va plus de quatre fois plus
vite dans l'eau que dans l'air.

Le reste de ces pages — les échos des Tuileries et du jardin de Descartes, la
lumière comme mouvement de la matière subtile, les renvois à la
\emph{Dioptrique} et aux \emph{Météores}, à Hérigone, à la « lettre de
M. Deschamps sur cette réfraction » — n'appelle pas de calcul.

\section*{Le livre du mouvement des corps}

\pagerange{11}{16}
\subsection*{Un corps qui tombe vers le centre d'une Terre qui tourne}

La plus longue remarque du volume commente la page 94 du livre. Qu'un corps
pesant tombe d'un point A de l'équateur vers le centre E de la Terre pendant
que la Terre tourne de son mouvement journalier. Galilée, dans son
\emph{Dialogue}, avait avancé que le corps décrirait un demi-cercle ; la
remarque soutient qu'il décrit une spirale, et en donne deux propriétés.

Les hypothèses sont celles du feuillet : le rayon EK qui porte le corps tourne
uniformément avec la Terre, et la distance tombée le long de ce rayon suit la
loi de Galilée, proportionnelle au carré du temps, jusqu'au centre. Si $R$ est
le rayon, $\Theta$ l'angle dont la Terre tourne pendant toute la chute et
$\theta$ l'angle déjà parcouru, la distance tombée est $R(\theta/\Theta)^2$ et
la courbe a pour équation polaire
\[
r(\theta) = R\left(1 - \frac{\theta^2}{\Theta^2}\right), \qquad
0 \le \theta \le \Theta .
\]
La première propriété du feuillet — « EK est à K21 en raison doublée de la
circonférence AKC à la circonférence AK » — est exactement cette loi :
$R : (R - r) = (\Theta : \theta)^2$.\footnote{« Raison doublée » : le carré du
rapport. Le point K21 est le point de la spirale sur le rayon EK ; la
transcription note que plusieurs fins de ligne de ce passage sont cachées dans
la reliure, sans que le sens en souffre.}

La seconde est que l'aire comprise entre la spirale et le rayon EA est au
secteur AKCE « comme 8 à 15 ». Elle est juste :
\[
\frac12\int_0^\Theta r(\theta)^2\,d\theta
 = \frac{R^2\Theta}{2}\int_0^1 (1-u^2)^2\,du
 = \frac{R^2\Theta}{2}\left(1 - \frac23 + \frac15\right)
 = \frac{8}{15}\cdot\frac{R^2\Theta}{2}.
\]
Le feuillet, qui n'a pas le calcul intégral, la démontre comme Archimède
démontrait l'aire de sa spirale. On divise l'arc en $n$ parties égales ($n = 16$
dans la figure) ; sur chaque partie, la spirale est prise entre un secteur
circulaire inscrit et un secteur circonscrit, de rayons
$R(1 - k^2/n^2)$ et $R(1 - (k-1)^2/n^2)$. Les secteurs étant semblables, leurs
aires sont comme les carrés de leurs rayons, et tout revient à comparer, en
prenant $R = n^2$,
\[
\sum_{k=1}^{n}(n^2 - k^2)^2 \quad\text{et}\quad \sum_{k=0}^{n-1}(n^2 - k^2)^2
\qquad\text{à}\qquad \frac{8}{15}\,n \cdot (n^2)^2 .
\]
C'est le « lemme » dont la remarque dit que « toute la difficulté consiste à le
prouver », et qu'elle se contente de vérifier pour $n = 16$ dans une table de
marge. La table est entièrement juste : les nombres $255, 252, 247, \dots, 31$
sont $256 - k^2$, leurs carrés $65025, 63504, \dots, 961$ ; la somme des quinze
carrés est $526472$, et $592008$ avec le carré $65536$ de $256$ ; le carré
$65536$ pris seize fois est $1048576$ ; et les produits croisés
$592008 \cdot 15 = 8880120 > 8388608 = 1048576 \cdot 8$ et
$526472 \cdot 15 = 7897080 < 8388608$ établissent que la somme circonscrite est
au-dessus et la somme inscrite au-dessous de $\frac{8}{15}$.\footnote{Le texte
écrit que « le premier secteur est au second comme le carré de 961 à celui de
3600 » : 961 et 3600 sont déjà les carrés de 31 et de 60, et il faut lire
« comme 961 à 3600 ». La transcription le signale.}

Le lemme est vrai pour tout $n$, et voici ce qui manque au feuillet pour le
démontrer. On a, en développant et en sommant les puissances,
\[
\sum_{k=1}^{n}(n^2-k^2)^2 = n^5 - 2n^2\sum_{k=1}^{n} k^2 + \sum_{k=1}^{n}k^4
 = \frac{8n^5}{15} - \frac{n^4}{2} - \frac{n}{30},
\]
qui est strictement inférieur à $\frac{8}{15}n^5$ ; et la somme circonscrite
vaut la même expression augmentée de $n^4$ (le terme $k = 0$), soit
$\frac{8n^5}{15} + \frac{n^4}{2} - \frac{n}{30}$, strictement supérieure. Les
deux sommes, divisées par $n^5$, diffèrent de $\frac1n$ et encadrent
$\frac{8}{15}$ : c'est la double inégalité qui fait marcher l'exhaustion, et
la remarque avait raison de dire qu'on trouverait le lemme « véritable en en
cherchant la démonstration ». Elle avait raison aussi d'ajouter que la
démonstration reste la même si l'angle $\Theta$ dépasse un demi-cercle, ou
fait plusieurs tours : rien dans le calcul ne dépend de $\Theta$. Et si le
corps part d'un parallèle et non de l'équateur, la courbe est tracée sur un
cône de sommet le centre de la Terre ; le cône se développe sur un plan sans
changer les aires, la courbe développée est une spirale de la même loi, et le
rapport $\frac{8}{15}$ demeure.\footnote{La page suivante renvoie à « l'hélice
décrite par M. Fermat touchant ladite descente, laquelle est ci-devant démontrée
sur la page 94 et 95 » : la copie rattache donc elle-même cette spirale à
Fermat. Elle rapporte aussi que Galilée « a depuis confessé par lettre qu'il
n'avait parlé de ce demi-cercle que par caprice et galanterie, et non à bon
escient ». Cette lecture n'a vérifié ni l'une ni l'autre de ces deux
affirmations sur des sources ; elle les rapporte comme la copie les donne.}

Il faut dire enfin que le modèle est celui du feuillet et non celui de la
mécanique. Dans la mécanique de Newton, à l'intérieur d'une Terre homogène la
pesanteur est proportionnelle à la distance au centre et non constante ; un
corps lâché à l'équateur avec la vitesse de rotation du sol décrit, dans un
repère fixe, une ellipse centrée au centre de la Terre, passe à quelque 375 km
du centre sans l'atteindre, et y arrive au bout d'environ 21 minutes, non de
plusieurs heures.\footnote{Ces nombres supposent une Terre homogène, sans
frottement, et un tunnel qui laisse le corps libre. La pulsation du mouvement
est $\Omega = \sqrt{g/R}$, d'où un quart de période de $\frac{\pi}{2}\sqrt{R/g}
\approx 21$ min et un demi-petit axe $\omega R/\Omega \approx 375$ km, $\omega$
étant la vitesse angulaire de la Terre. Une remarque de la page 96 précise que
le « il devrait tomber en 6 heures » du livre s'entend de la surface jusqu'au
centre.} La spirale des $\frac{8}{15}$ est donc un théorème exact sur une
hypothèse fausse.

\subsection*{La courbe n'est pas un cercle}

Une seconde démonstration, en algèbre (vue 13), établit que la courbe ne peut
être le demi-cercle de Galilée. Soit A le centre de la Terre, $AC = b$ le rayon,
et le cercle de diamètre AC que Galilée proposait. Si le corps est en K sur ce
cercle quand la Terre a tourné de l'angle CAH, et en I quand elle a tourné
du double, CAD, la loi de chute veut que la chute ID soit quadruple de la chute
HK. Le feuillet pose $AK = e$, calcule $ID = \dfrac{b^2 - e^2}{\frac12 b}$ et
$HK = b - e$, et montre que $ID = 4HK$ équivaudrait à $2be = b^2 + e^2$,
c'est-à-dire à $(b - e)^2 = 0$, ce qui n'est pas. En notation moderne : si
l'angle KAC est $\varphi$, $e = b\cos\varphi$ et $AI = b\cos2\varphi$, d'où
\[
\frac{ID}{HK} = \frac{b(1-\cos2\varphi)}{b(1-\cos\varphi)}
 = 2(1+\cos\varphi) = \frac{2(b+e)}{b},
\]
qui est toujours compris strictement entre 2 et 4 et n'atteint 4 qu'à la
limite $e = b$. Le calcul du feuillet est juste, et sa conclusion aussi : le
cercle ne satisfait la loi de chute en aucun point.\footnote{Dans
l'expression de $AI = CL$, le feuillet écrit un $+$ et un $-$ superposés,
« $\dfrac{+ee \mp \frac12 bb}{\frac12 b}$ » ; le calcul du rectangle LAE,
$\frac12 ab + \frac12 b^2 = e^2$, donne sans ambiguïté le signe $-$ :
$AI = \dfrac{e^2 - \frac12 b^2}{\frac12 b} = b\cos2\varphi$. Le feuillet emploie
la double barre « $\parallel$ » pour l'égalité et pour « est à \dots comme ».}

La suite montre davantage : pour tout rapport $\lambda$ compris entre 2 et 4,
il existe un point K et un seul tel que $ID : HK = \lambda$, à savoir celui où
$AK = e = (\frac{\lambda}{2} - 1)b$. Le feuillet l'écrit avec
$\lambda = s : \frac12 b$, et trouve $e = s - b$ : c'est la même chose. Il
conclut que la courbe de chute est une droite sous le pôle, une spirale plane
sous l'équateur et, ailleurs, une spirale tracée sur un cône isocèle de sommet
le centre de la Terre — ce qui est juste dans le modèle.

\subsection*{Les spirales de tout degré, et une table sans titre}

La remarque sur la page 98 pose le problème général : si la chute est
proportionnelle à une puissance quelconque du temps, trouver le rapport de
l'aire de la spirale au cercle. « La première hélice est d'Archimède, et la
seconde celle de Galilée pour les graves ; les autres ne tombent point en la
pratique de la physique. » La réponse est simple. Si la distance tombée est
$R(\theta/\Theta)^n$, l'aire est au secteur comme
\[
\int_0^1 (1-u^n)^2\,du = 1 - \frac{2}{n+1} + \frac{1}{2n+1}
 = \frac{2n^2}{(n+1)(2n+1)} .
\]
Pour $n = 1, 2, \dots, 8$ ce rapport vaut $\frac13$, $\frac{8}{15}$,
$\frac{9}{14}$, $\frac{32}{45}$, $\frac{25}{33}$, $\frac{72}{91}$,
$\frac{49}{60}$, $\frac{128}{153}$. Or la marge gauche de la même page porte, sans
titre et sans renvoi, une petite table de deux colonnes : $3 \mid 1$ ;
$15 \mid 8$ ; $14 \mid 9$ ; $45 \mid 32$ ; $33 \mid 25$ ; $91 \mid 72$ ;
$60 \mid 49$ ; $153 \mid 128$. Ce sont, dénominateur puis numérateur et réduits
au plus bas terme, exactement ces huit rapports. La table est la solution du
problème de la page 98 pour les huit premiers degrés ; le premier est le
tiers qu'Archimède avait trouvé pour sa spirale, le second les $\frac{8}{15}$
de la page 94.\footnote{L'identification de la table est de cette lecture ;
la transcription la décrit sans l'interpréter, et la copie ne l'appelle
nulle part. L'énoncé de la page 98, lu à la lettre, fait la chute
proportionnelle à une racine du temps et non à une puissance : « même
proportion de la circonférence à l'arc que de AB à FC, ou que du carré de AB au
carré de FC ». Avec cette lecture on trouverait $\frac{2}{(n+1)(n+2)}$, soit
$\frac16$ pour la spirale de Galilée, en contradiction avec les $\frac{8}{15}$
démontrés deux pages plus haut. C'est la table qui est cohérente avec la page 94,
et c'est l'énoncé qui a les puissances du mauvais côté.}

\subsection*{La loi des nombres impairs}

La remarque sur la page 100 explique d'où viennent les tables de chute. Si le
mobile parcourt AB dans la première minute, la vitesse acquise en B lui ferait
parcourir $2AB$ dans une minute égale ; la pesanteur ajoutant encore AB, il
parcourt $3AB$ dans la deuxième minute. En C, parvenu en deux minutes à
$AC = 4AB$, sa vitesse lui ferait faire $2AC = 8AB$ en deux minutes, soit $4AB$
par minute, et la pesanteur en ajoute un : $5AB$. En D, $9AB$ en trois minutes,
$6AB$ par minute, plus un : $7AB$. Les espaces parcourus en temps égaux croissent
comme les nombres impairs $1$, $3$, $5$, $7$, et la marge dessine l'échelle. Le
raisonnement est juste : il repose sur la propriété, établie par Galilée, que
la vitesse acquise en un temps $t$ ferait parcourir d'un mouvement uniforme le
double de l'espace tombé pendant ce temps ; c'est la relation
$d(t) = \frac12 g t^2$, dont les accroissements sont
$d(k) - d(k-1) = (2k-1)\,d(1)$.

Les remarques voisines disent que la cycloïde « enferme un espace triple du
cercle qui la décrit » (vrai : l'aire sous une arche est $3\pi r^2$), que la
courbe de la page 120 n'est pas une demi-ellipse mais cette cycloïde, et, à la
page 101, que si les graves descendent par le mouvement de la matière subtile
« le fondement précédent ne sert pas ».\footnote{« Trochoïde » et « roulette »
sont les noms du temps pour la cycloïde.}

\subsection*{La résistance des milieux}

Les remarques des pages 129 à 131 examinent la règle de Galilée selon laquelle un
milieu retranche à la vitesse d'un corps la fraction de son poids qu'il lui
ôte : l'eau ôterait le quart de la vitesse d'une pierre quatre fois plus
pesante qu'elle, et le douzième de celle du plomb. Mille brasses faites en 24
secondes dans le vide le seraient, dit la remarque, en 30 secondes par la pierre
et en 26 par le plomb. Ces deux nombres allongent le temps du quart et du
douzième ($24 + 6$, $24 + 2$) ; si c'est la vitesse qui est réduite dans ces
proportions, les temps sont $32$ et $26\frac{2}{11}$ secondes.\footnote{De même
les « points d'égalité » à 100 et 90 brasses, qui supposent des vitesses
limites dans le rapport $10 : 9$, ne suivent pas de la règle, qui donnerait
$\frac{11}{12} : \frac34 = 11 : 9$. La copie note elle-même que « tout ce
discours de Galilée a grand besoin d'examen ».} La remarque de la page 131 —
les graves n'augmentent pas indéfiniment leur vitesse, mais approchent d'une
limite « suivant les approchements des asymptotes » — énonce ce qu'on appelle
aujourd'hui la vitesse limite ; elle est juste, et c'est la résistance
de l'air, qui croît avec la vitesse, qui en est la cause. Les rapports donnés
pour la moelle de sureau et le plomb varient d'une remarque à l'autre (360 dans
le livre, 193 page 129, $128\frac23$ page 142).

La remarque sur la page 135 corrige la longueur du pendule qui bat la seconde :
le sieur Decousu dit deux pieds dix pouces, la copie répond que « trois pieds
suffiront ». La longueur vraie, pour un pendule dont chaque oscillation simple
dure une seconde, est de $0{,}994$ m, soit trois pieds et huit lignes et
demie : la copie est plus près du vrai que Decousu.

\subsection*{Les pendules composés}

Les remarques des pages 136 et suivantes comparent au « filet isochrone » — la
longueur du pendule simple qui bat en même temps — toutes sortes de corps
suspendus : bâton, triangles, cercle, sphère, parabole, cycloïde. La théorie
moderne de ces pendules est celle du pendule pesant : un corps de masse $m$,
de moment d'inertie $I$ par rapport à l'axe de suspension, dont le centre de
gravité est à la distance $d$ de l'axe, oscille comme un pendule simple de
longueur $\ell = I/(md)$. Elle permet de vérifier plusieurs articles.

Le bâton suspendu par un bout a une longueur égale aux $\frac32$ du filet
isochrone : la copie dit « en raison sesquialtère proxime », et c'est
exactement $\frac32$ ($\ell = \frac23 L$). La variante qu'elle rapporte du
P. Fabri, « comme l'aire du demi-cercle au carré de son rayon », soit
$\frac{\pi}{2} \approx 1{,}57$, est fausse.\footnote{« Sesquialtère » : le
rapport $\frac32$. « Proxime » : à peu près.}

Pour une plaque triangulaire isocèle de hauteur $h$ et d'angle au sommet
$\alpha$, oscillant dans son plan, on trouve $\ell = \frac34 h +
\frac{h}{4}\tan^2\frac{\alpha}{2}$ si elle est suspendue par la pointe, et
$\ell = \dfrac{h}{2\cos^2(\alpha/2)}$ si elle est suspendue par le milieu de
la base. Rapportées à la hauteur, ces valeurs se comparent à celles de la copie
ainsi :
\[
\begin{array}{r|cc|cc}
\alpha & \text{pointe, théorie} & \text{pointe, copie}
  & \text{base, théorie} & \text{base, copie} \\ \hline
20^\circ & 1 : 0{,}76 & \text{presque } 4:3 & 1 : 0{,}52 & 3 : 2 \\
60^\circ & 6 : 5 & 6 : 5 & 3 : 2 & 3 : 2 \\
90^\circ & 1 : 1 & 1 : 1 & 1 : 1 & 1 : 1 \\
110^\circ & 1 : 1{,}26 & 2 : 3 & 1 : 1{,}52 & 8 : 15 \\
120^\circ & 2 : 3 & \text{« quasi de même »} & 1 : 2 & \text{« quasi de même »} \\
135^\circ & 1 : 2{,}21 & 2 : 5 & 1 : 3{,}41 & 1 : 4 \\
150^\circ & 1 : 4{,}23 & 1 : 4 & 1 : 7{,}46 & 1 : 12 \\
170^\circ & 1 : 33{,}4 & 1 : 31 & 1 : 65{,}8 & 1 : 36
\end{array}
\]
L'accord est exact pour le triangle de $60^\circ$ par la pointe ($6 : 5$) et par
la base ($3 : 2$), et pour le triangle rectangle, « égal au filet » des deux
façons ; il est bon, pour la pointe, jusqu'à $90^\circ$ et à $150^\circ$. Pour
les triangles obtus suspendus par la base, les valeurs de la copie s'écartent
de la théorie, et le « $2 : 3$ » donné pour $110^\circ$ par la pointe est la
valeur exacte de $120^\circ$.\footnote{Cette lecture suppose que la copie
compare au filet la hauteur du triangle, et que « par la base » veut dire
« par le milieu de la base » ; c'est l'hypothèse sous laquelle les trois
valeurs exactes coïncident. La copie ne le précise pas, et ses nombres, qui
sont des mesures, sont donnés comme approchés.}

Parmi les autres corps : la circonférence (un anneau) suspendue par un point
est « égale » au filet, et c'est juste, $\ell$ étant le diamètre ; le cercle
plein suspendu par un point de son bord a un diamètre qui est au filet comme
$4 : 3$ et non « comme 7 à 5 » ; la sphère suspendue par un point de sa surface
est au filet comme $10 : 7 \approx 1{,}43$, et la copie dit « comme $1\frac49$ à
un » ($1{,}44$), ce qui est bon. Le demi-cercle, par le convexe ou par la base,
donne en théorie $1 : 1{,}13$ et $1 : 1{,}18$ du rayon au filet ; la copie dit
« comme 5 à 6 » pour les deux. Les autres articles (quart de cercle, cône,
demi-sphère, parabole, roulette, « parabole seconde de Torricelli ») ne
disent pas assez comment le corps est suspendu pour être vérifiés.

\section*{Le troisième livre du mouvement : cordes, poids, rupture, projectiles}

\pagerange{16}{19}
\subsection*{L'ellipse et le triangle}

La remarque sur la page 157 compare deux familles de triangles de même base AB.
Si le sommet décrit l'ellipse de foyers A et B, tous les triangles ont le même
périmètre, et le triangle isocèle a la plus grande aire. Si le sommet décrit
une parallèle à AB, tous ont la même aire, et le triangle isocèle a le plus
petit périmètre. Les deux énoncés sont justes et duaux ; mais la copie conclut
que « tant entre parallèles que dans l'ellipse, le triangle isopleure est le
plus petit \emph{in ambitu} », ce qui est faux pour l'ellipse, où les périmètres
sont égaux.\footnote{« Isopleure », ici, veut dire isocèle, puisque le
triangle a pour base la corde donnée. « In ambitu » : par le pourtour.}

\subsection*{La corde qui revient, et la loi de la tension}

La remarque sur la page 160 découpe le retour d'une corde tendue en 88 parties
de temps : $8$, $20$, $28$, $32$ espaces dans les quatre temps successifs, dont la
somme fait bien 88. C'est un modèle, non une démonstration.

La table de marge de la vue 18 donne les poids qui, tendant huit cordes égales,
leur font sonner les huit degrés de l'octave, de 6 livres pour l'ut grave à 24
pour l'ut aigu. La loi est que la fréquence est proportionnelle à la racine
carrée de la tension : le poids est donc $6 \cdot f^2$ si $f$ est le rapport de
fréquence à l'ut grave. Avec les rapports de l'intonation juste, on trouve
$16\frac23$ pour la ($f = \frac53$), $13\frac12$ pour sol, $10\frac23$ pour fa,
$7\frac{33}{81}$ pour ré si l'on prend le ton mineur $\frac{10}{9}$ — ce que
fait la table — et 24 pour l'ut aigu : ces valeurs de la table sont justes.
Deux ne le sont pas : pour mi ($f = \frac54$), $6 \cdot \frac{25}{16} =
9\frac38$ et non $9\frac18$ ; pour si ($f = \frac{15}{8}$), $6\cdot
\frac{225}{64} = 21\frac{3}{32}$ et non $21\frac{3}{16}$.

La remarque sur la page 189 justifie par un raisonnement sur des triangles
que la corde à l'octave doit être tendue quatre fois plus. La conclusion est
juste, et c'est ce qu'on appelle aujourd'hui l'une des lois de Mersenne : la
fréquence fondamentale d'une corde de longueur $L$, de masse linéique $\mu$,
tendue par la force $T$, est $f = \frac{1}{2L}\sqrt{T/\mu}$, de sorte que
doubler $f$ demande de quadrupler $T$. Le raisonnement de la copie, qui égale
entre les deux cordes le produit de la force par le temps, n'est pas une
démonstration.\footnote{Le nom de « lois de Mersenne » est un usage moderne ;
la copie ne nomme aucune loi. La remarque sur la page 169 rapporte qu'une corde
à l'unisson d'un tuyau d'orgue bouché de 2 pieds bat l'air 108 fois par
seconde ; la transcription note que « 108 » peut aussi se lire « 128 ». Un tuyau
bouché de 2 pieds de roi ($0{,}65$ m) donne, par la formule élémentaire
$f = c/4L$, environ 130 vibrations par seconde : la lecture 128 est la plus
proche.}

\subsection*{Un nombre « qui semble faux » et qui est juste}

La remarque sur la page 193 rapporte la règle de Galilée : la résistance des
cylindres à la rupture est en raison triplée des diamètres, de sorte qu'un
bâton deux fois plus mince est huit fois plus aisé à rompre. Puis elle calcule :
une corde de laiton de $\frac16$ de ligne de diamètre rompant sous 18 livres,
une colonne d'un pied de diamètre romprait « seulement avec 13436928 » — et la
marge ajoute : « ce nombre me semble faux ».

Le rapport des diamètres est $144 : \frac16 = 864$, un pied valant 144
lignes. Or $18 \cdot 864^2 = 13\,436\,928$ exactement : le nombre suit la loi
des carrés des diamètres, c'est-à-dire des aires des sections, et non celle des
cubes, qui donnerait $18 \cdot 864^3 = 11\,609\,505\,792$. C'est la loi des
carrés qui est la bonne ici. La règle des cubes de Galilée vaut pour une
poutre rompue par flexion ; une corde qui casse sous un poids suspendu est
rompue par traction, et sa résistance est proportionnelle à l'aire de sa
section. Le nombre est juste pour le problème posé ; c'est la phrase qui le
précède qui ne s'y applique pas.

\subsection*{Chute et projectiles}

La remarque sur la page 197 corrige le livre : à la fin du quatrième pied de
chute, le mobile n'a que deux degrés de vitesse, puisque la vitesse croît comme
la racine de l'espace ($\sqrt4 = 2$), et cette vitesse, gardée uniforme, lui
ferait parcourir huit pieds dans le temps qu'il en a mis à tomber de quatre.
C'est juste. La remarque sur la page 222 dit qu'une flèche qui met 12 secondes à
monter et autant à descendre ne monte que de 288 toises : c'est la hauteur
$12 \cdot 12^2 = 1728$ pieds, avec la valeur de 12 pieds pour la chute de la
première seconde que le volume emploie partout. La valeur vraie de cette
chute, sans air, est de $4{,}9$ m, un peu plus de 15 pieds. Une remarque
ajoute que d'autres expériences montrent la flèche plus prompte à monter qu'à
descendre — trois secondes contre cinq — ce qui est l'effet de la résistance de
l'air et contredit la symétrie supposée par le livre. La remarque sur la page
173 (quatre fois plus de force pour envoyer les missiles deux fois plus loin)
est juste dans le vide, où la portée est proportionnelle au carré de la
vitesse initiale.

\section*{Le livre de la voix}

\pagerange{19}{19}
Trois remarques. Les huit poids « en progression triple » sont
$1$, $3$, $9$, $27$, $81$, $243$, $729$, $2187$, les puissances $3^0$ à $3^7$. Une remarque de
métaphysique, sur la page 80, conteste que ce qui est indépendant doive avoir
toute perfection infinie. Enfin, à la ligne 38 de la même page : « M. de
Villiers a observé l'année 1638 qu'au crâne d'un sourd et muet, des trois
osselets qui se trouvent en l'oreille, l'incus, le malleolus et le stapes, le
premier manquait en chaque oreille. » C'est l'une des dates écrites dans le
volume.\footnote{La transcription note que « 1638 » a été lu à fort
grossissement et que « stapes » est lu au bord de la reliure.}

\section*{Le livre des chants : permutations et combinaisons}

\pagerange{19}{24}
\subsection*{Les mots de vingt-trois lettres}

La remarque sur la page 107 renvoie au « problème arithmétique » de Guldin :
combien de mots peut-on faire avec les 23 lettres de l'alphabet, aucune lettre
n'étant répétée dans un mot, et combien de lettres tous ces mots contiennent-ils.
Le nombre de mots de $k$ lettres est $23!/(23-k)!$ ; le nombre total de mots est
\[
W = \sum_{k=1}^{23}\frac{23!}{(23-k)!} = \lfloor e \cdot 23! \rfloor - 1
 = 70\,273\,067\,330\,330\,098\,091\,155,
\]
et le nombre total de lettres
$L = \sum_k k\cdot\frac{23!}{(23-k)!} = 1\,546\,007\,481\,267\,262\,158\,005\,433$.
La copie donne pour les lettres $1546007491267202147905433$ et pour les mots
$702730673303300980911155$. Le premier nombre a le bon ordre de grandeur et les
huit premiers chiffres justes, mais diffère au milieu ; le second a un chiffre de
trop — un « 1 » de plus avant la fin — alors que la copie dit elle-même que son
premier chiffre « vaut septante sextillions », c'est-à-dire que le nombre a 23
chiffres, comme $W$.\footnote{Dans la numération de la copie, le sextillion est
$10^{21}$ et le septillion $10^{24}$. Les deux nombres sont lus chiffre à chiffre
par la transcription ; cette lecture ne peut dire si l'écart vient de Guldin,
de Mersenne, de la copie ou de la lecture. La formule $\sum_{k=0}^{n} n!/k! =
\lfloor e\,n! \rfloor$ est moderne.}

La remarque sur la page 110 dit que le 64\textsuperscript{e} nombre de la table du
livre a 90 chiffres. Si ce nombre est $64!$, c'est juste : $64!$ a 90 chiffres.
Mais il commence par 126, et non par les « 221 » que la remarque donne pour son
« premier ternaire » ; sans la table du livre, cette lecture ne peut trancher.

\subsection*{Ranger toutes les permutations}

La remarque sur la page 111 donne la règle pour écrire toutes les permutations
de DIEU sans en oublier ni en répéter : on garde la première lettre et l'on
permute les trois autres, en gardant à son tour la première de celles-ci, et
ainsi de suite ; puis on met la deuxième lettre en tête, et l'on recommence.
C'est l'énumération dans l'ordre lexicographique, et elle démontre en même
temps la récurrence $n! = n \cdot (n-1)!$ : $2 \cdot 3 = 6$, $6 \cdot 4 = 24$,
$24 \cdot 5 = 120$, $120 \cdot 6 = 720$, $720 \cdot 7 = 5040$. Tout est
juste.\footnote{« La combinaison de l'ordre » de $n$ choses est, dans la langue
de la copie, le nombre de leurs arrangements, $n!$.}

\subsection*{Les chants de huit notes}

La remarque sur la page 113 retire des chants de l'octave ceux qui contiennent
une septième majeure ou une sixte. Il y a $8! = 40320$ chants de huit notes
différentes.\footnote{La copie écrit « des 4320 chants de l'octave » ; toute la
suite du calcul suppose 40320, et le chiffre manquant est une faute de copie ou
de plume.} Ceux où ut et si sont voisins se comptent en collant ut et si en une
seule note : $7! = 5040$, doublé pour les deux ordres, $10080$ ; de même pour ut
et la ; en tout $20160$. Mais ceux où ut est voisin à la fois de si et de la
ont été comptés deux fois : on colle « si ut la » ou « la ut si » en une note,
$2 \cdot 6! = 1440$, et l'on retire. Il reste $18720$ chants qui contiennent l'un
ou l'autre intervalle, et donc $40320 - 18720 = 21600$ qui n'en contiennent
aucun. C'est exactement le principe d'inclusion et d'exclusion, correctement
appliqué.

\subsection*{Le piquet}

La remarque sur la page 131 distingue deux sortes de « combinaisons ». Quand
toutes les choses sont différentes, comme les douze cartes d'une main prises
dans les 36 du jeu de piquet, le nombre de mains ordonnées est
$36 \cdot 35 \cdots 25 = 599\,555\,620\,984\,320\,000$, et le nombre de mains,
l'ordre ne changeant pas le jeu, est ce produit divisé par $12! = 479\,001\,600$,
soit $\binom{36}{12} = 1\,251\,677\,700$. Les deux nombres sont justes.\footnote{À
la vue 23, la copie écrit le même nombre « 125677700 », avec un chiffre de
moins ; à la vue 22 il est juste.} Quand les choses peuvent être semblables,
comme douze cartes prises dans douze jeux, il faut multiplier $36 \cdot 37
\cdots 47$ et diviser par $12!$ : c'est le nombre des combinaisons avec
répétition, $\binom{47}{12} = 52\,251\,400\,851$. De même, cent fruits choisis
dans six paniers donnent $6 \cdot 7 \cdots 105 / 100! = \binom{105}{100} =
\binom{105}{5} = 96\,560\,646$ choix.\footnote{La copie écrit « multiplier 37 par
37 » ; il faut 36 par 37, comme le montre le dernier facteur, 47. La règle est
dite plus loin (vue 23) « de M. Fren. », abréviation que la copie ne développe
pas.}

\subsection*{Les « puissances triangulaires »}

Ce que la copie appelle la $k$-ième « puissance triangulaire » de la « racine »
$r$ est le nombre figuré $\binom{r+k-1}{k}$ : la deuxième est le nombre
triangulaire, la troisième le tétraédrique, et chacune est la somme des $r$
premières de la précédente. La règle de calcul qu'elle donne — multiplier
$r, r+1, \dots, r+k-1$ et diviser par $k!$ — est exacte, et l'exemple aussi : la
sixième puissance triangulaire de 5 est $5 \cdot 6 \cdot 7 \cdot 8 \cdot 9
\cdot 10 / 720 = 210 = \binom{10}{6}$, et l'astuce qui évite la division
(retirer $1, 2, \dots, 6$ des facteurs, il reste $7 \cdot 3 \cdot 10$) est juste.

Mais la copie écrit que les douze cartes sans ordre sont « la
12\textsuperscript{e} puissance triangulaire de 36, qui est à la fin de ladite
table, à savoir 1251677700 ». Le nombre est juste, le nom ne l'est pas. Dans sa
propre définition, la 12\textsuperscript{e} puissance triangulaire de 36 est
$\binom{47}{12}$, le nombre des mains avec répétition ; et $1251677700 =
\binom{36}{12}$ est la 12\textsuperscript{e} puissance triangulaire de 25 —
précisément le dernier nombre d'une table qui, dit la copie, va « jusqu'à la
12\textsuperscript{e} puissance de 25 ».

La « combinaison générale » de cent choses de six sortes, l'ordre compris, est
$6^{100}$, et la copie a raison de dire qu'on ne peut en tirer le nombre sans
ordre par une seule division : sur l'exemple de quatre choses, les types abcd,
aabc, aaab, aabb, aaaa se permutent en $24$, $12$, $4$, $6$, $1$ façons, de somme 47,
nombre premier, et chaque type demande son propre diviseur.

\subsection*{Une correction qui corrige mal}

La remarque sur la page 139 reprend un calcul du livre : il faut multiplier
« 484, qui est le troisième nombre de la table », par 15, ce qui donne 7260 ;
« 10648 », le quatrième nombre de la cinquième colonne, par 16 ; et « le
nombre cinquième, 234276 », par 15, « dont il viendra 3514140, et non 3513840,
comme il y a dans le texte qui est fautif ». Les produits sont justes. Mais
$484 = 22^2$ et $10648 = 22^3$, et le nombre suivant de la colonne doit être
$22^4 = 234256$ ; or $234256 \cdot 15 = 3513840$, le nombre même du texte
imprimé. Si la colonne est bien celle des puissances de 22, c'est le livre qui
avait raison, et la correction porte sur un 234276 fautif.\footnote{Cette lecture
n'a pas ouvert le livre et ne peut dire si le 234276 est une coquille de la table
imprimée, une faute de la copie ou une lecture ; elle constate seulement que
3513840 est $15 \cdot 22^4$.}

\subsection*{Petits corollaires}

La somme des nombres de 1 à $n$ est donnée par la règle de Gauss avant la
lettre : $\frac{n}{2}(n+1)$ si $n$ est pair ($7 \cdot 15 = 105$ pour 14),
$n \cdot \frac{n+1}{2}$ s'il est impair ($5 \cdot 3 = 15$, $15 \cdot 8 = 120$).
Deux choses prises dans dix, avec répétition et ordre, font $10^2 = 100$
variétés ; trois dans neuf, $9^3 = 729$. Le produit de 36 par 35, divisé par 2,
donne $\binom{36}{2} = 630$. Tout est juste.

\subsection*{Les chants à quatre parties}

La remarque sur la proposition 20 compte $2\,520\,473\,760\,000$ chants à quatre
parties. Ce nombre est $1260^4$, et il faut, à un chant par seconde, un peu moins
de 80\,000 ans pour les chanter tous, soit 800 siècles : l'affirmation qu'« il
faudrait plus de cent mille siècles » ne suit pas, non plus que le nombre de
45\,048\,300 chants en cent ans, qui ne correspond pas à un chant par seconde
(un siècle compte environ $3{,}16 \cdot 10^9$ secondes). Pour quatre notes
permutées dans chaque partie, les duos sont $24^2 = 576$, ce que dit la copie ;
les trios $24^3 = 13824$ et les quatuors $24^4 = 331776$, et non 13729 et
319376.\footnote{La transcription relève elle-même que 13729 et 319376 « ne sont
pas les puissances de 24 auxquelles 576 ferait attendre ».}

\subsection*{Les carrés magiques}

La même page 110 annonce que le nombre des manières de varier un carré magique
en le gardant magique est plus grand encore, et donne pour le carré de 22 un
produit de trois nombres de 34, 32 et 33 chiffres, pour le carré de 14 le
nombre 33710564646400. La copie redonne les mêmes nombres sur les feuillets de la
fin (vue 46), et les deux lectures ne s'accordent pas : les deuxième et
troisième facteurs y ont un zéro final de moins, et le nombre du carré de 14 y
est lu 3371056466400, suivi d'un chiffre taché. Cette lecture ne sait pas
reconstituer ces dénombrements et ne les vérifie pas.\footnote{Les deuxième et
troisième facteurs n'ont que des diviseurs premiers au plus égaux à 37, comme un
produit de factorielles ; le premier a un grand facteur premier et ne peut en
être un.}

\section*{Les livres des consonances}

\pagerange{24}{28}
\subsection*{Diviser l'octave et la quarte en trois raisons superparticulières}

Une raison superparticulière est un rapport de la forme $\frac{n+1}{n}$.
L'octave, dit la remarque sur la page 55, ne se divise en trois raisons
superparticulières « qu'en cinq sortes », et la marge les donne :
$\frac32\cdot\frac54\cdot\frac{16}{15}$, $\frac32\cdot\frac65\cdot\frac{10}{9}$,
$\frac32\cdot\frac76\cdot\frac87$, $\frac43\cdot\frac43\cdot\frac98$,
$\frac43\cdot\frac54\cdot\frac65$. C'est exact et complet : l'équation
$\bigl(1+\frac1a\bigr)\bigl(1+\frac1b\bigr)\bigl(1+\frac1c\bigr) = 2$ avec
$a \le b \le c$ entiers force $a \le 3$ (sinon le produit est au plus
$(\frac54)^3 < 2$), puis borne $b$, puis détermine $c$, et une énumération
complète donne ces cinq solutions et elles seules.

La quarte, dit la remarque sur la page 142, « se peut diviser 27 fois en trois
raisons superparticulières, bien que Bryennios n'en tienne que 15 possibles »,
et la copie dresse la table des 27 divisions. Chacune des 27 est juste : le
produit de ses trois raisons vaut $\frac43$. Mais la même énumération complète,
faite sur $\bigl(1+\frac1a\bigr)\bigl(1+\frac1b\bigr)\bigl(1+\frac1c\bigr) =
\frac43$, ne donne que 26 divisions distinctes, et la table les contient toutes :
les divisions 14 et 24 de la table,
$\frac{16}{15}\cdot\frac76\cdot\frac{15}{14}$ et
$\frac76\cdot\frac{15}{14}\cdot\frac{16}{15}$, sont la même, écrite dans un autre
ordre. La quarte se divise en 26 manières et non en 27.\footnote{Si l'on comptait
comme différentes les divisions qui ne diffèrent que par l'ordre des raisons —
ce qui aurait un sens pour des tétracordes, où l'ordre des intervalles compte —
on en trouverait bien davantage que 27 ; la table ne le fait pour aucune autre
division. Les trois derniers termes des divisions 5, 10 et 15, cachés dans le pli
d'après la transcription, se complètent sans ambiguïté par le calcul :
$\frac{96}{95}$, $\frac{40}{39}$ et $\frac{100}{99}$.}

\subsection*{Tempérament et intonation juste}

La remarque sur la page 58 avertit qu'une proposition du livre n'est vraie
que « des vrais intervalles » et non du tempérament des instruments : sur un
instrument à tons égaux, trois quintes font exactement une treizième majeure
et trois tierces majeures font l'octave, « ce qui ne peut arriver en la vraie
théorie ». C'est exact : trois quintes justes font $\frac{27}{8}$, qui dépasse la
treizième majeure juste $\frac{10}{3}$ d'un comma syntonique $\frac{81}{80}$, et
trois tierces majeures justes font $\frac{125}{64}$, qui manque l'octave d'un
dièse $\frac{128}{125}$ ; dans le tempérament égal, $7 \cdot 3 = 21 = 12 + 9$ et
$4 \cdot 3 = 12$ demi-tons.

La remarque suivante somme des séries géométriques : les moitiés successives de
la moitié d'une corde font la moitié entière,
$\frac14 + \frac18 + \dots = \frac12$ ; les tiers successifs d'un tiers font la
moitié, $\frac13 + \frac19 + \dots = \frac12$ ; tous les quarts font
$\frac13$ et tous les dixièmes $\frac19$. C'est
$\sum_{k \ge 1} q^k = \frac{q}{1-q}$, correctement appliquée.

Les remarques des pages 71 à 76 reprennent la théorie des coïncidences : à
la quarte, trois battements de l'aiguë sur quatre sont désunis, à la quinte
deux sur trois, à l'octave un sur deux ; à la tierce majeure un seul sur
cinq s'unit. La remarque sur la page 91 donne les moindres termes des trois
médiétés, arithmétique $3, 2, 1$, géométrique $4, 2, 1$, harmonique $6, 4, 3$ :
tous justes ($4 = \frac{2\cdot 6\cdot 3}{6 + 3}$). La règle de trois de la
table 57 (page 109) est juste : $25 \cdot \frac89 = 22\frac29$, et
$4\cdot\frac{64}{81} = 3\frac{13}{81} = \bigl(1\frac79\bigr)^2$.

\subsection*{Les tables de Jean Le Fèvre}

La remarque sur la page 146 recopie les nombres dont Jean Le Fèvre se servait,
d'après Gafurius, pour les six voix. La table A donne les longueurs de la
gamme pythagoricienne, $27, 24, 21\frac13, 20\frac14, 18, 16$ (ut, ré, mi, fa,
sol, la), leurs carrés $729, 576, 455\frac19, 410\frac{1}{16}, 324, 256$ et
leurs cubes $19683, 13824, 9709\frac{1}{27}, 8303\frac{49}{64}, 5832, 4096$ :
tout est exact. La table C donne une autre gamme pythagoricienne, de
$10\frac18$ à 4, et les carrés de ses termes ; tout y est exact aussi, le
carré de $5\frac{1}{16}$ étant $25\frac58 + \frac{1}{256}$.\footnote{La
transcription lit la fraction isolée « $\frac{1}{1256}$ », le « 1256 » douteux :
il faut $\frac{1}{256}$.}

La table B est d'une autre nature que ne le dit le texte : ses poids, de 4608 à
10368 grains, ne sont pas les cubes des longueurs mais leur sont
proportionnels, l'octave doublant le poids ($4608 \to 9216$). Les degrés
diatoniques y sont justes, et leurs conversions en onces, gros et grains aussi
(une once valant 576 grains) ; trois des degrés chromatiques sont faux :
$5184 \cdot \frac{256}{243} = 5461\frac13$ et non $5461\frac18$,
$6912 \cdot \frac{256}{243} = 7281\frac79$ et non $7281\frac19$,
$6144 \cdot \frac{256}{243} = 6472\frac{56}{81}$, écrit 6472. Le total
annoncé, « $86681\frac{17}{27}$ grains, ou 9 livres 6 onces 3 gros 65 grains et
$\frac{17}{27}$ », est cohérent dans ses unités ; il n'est pas la somme des
quinze poids de la table ($106757\frac{19}{24}$), mais, à un grain près, celle
des treize premiers, d'un la à l'autre.

La remarque sur la page 151 ajoute que, selon l'ancien système d'Arezzo et de
Pythagore, la tierce majeure est $\frac{81}{64}$ et la mineure $\frac{32}{27}$,
ce qui est juste, et la « septième majeure » $\frac{27}{16}$, ce qui ne l'est
pas : $\frac{27}{16}$ est la sixte majeure pythagoricienne ; la septième majeure
est $\frac{243}{128}$.

Les remarques des pages 154 à 315 portent sur les modes, les espèces
d'octave et les « cadences » d'après Doni et le sieur Vincent ; elles ne
comportent pas de calcul, sinon que douze modes se réduisent à sept espèces,
selon la place des deux demi-tons dans l'octave, ce qui est juste.

\section*{Les livres des instruments}

\pagerange{28}{37}
\subsection*{Cordes}

La remarque sur la page 3 du livre des instruments reprend les recettes des
cordes de boyau déjà lues sur les feuillets de devant, et répond à la croyance
qu'une corde de loup ne s'accorde jamais avec une corde de mouton : il suffit
de les tendre pour qu'elles fassent le même nombre de vibrations en même temps,
« et partant elles feront l'unisson ». C'est la bonne réponse. La remarque sur
la page 5 énonce que des cylindres de même volume ont des surfaces latérales
« en raison sous-doublée de leurs hauteurs » : c'est exact, car pour un volume
$V$ fixé, $2\pi r h = 2\sqrt{\pi V h}$.\footnote{« Raison sous-doublée » : la
racine carrée du rapport.}

\subsection*{Le système de vingt-quatre quarts de ton}

La marge de la vue 29 donne un « système de 24 quarts de tons ou de 25 cordes, où
tout est égal » : les longueurs de 25 cordes, de 100000 à 50000, et une seconde
colonne de 200 à 100. C'est la division de l'octave en 24 parties égales, la
corde de rang $k$ valant $100000 \cdot 2^{-(k-1)/24}$. Les valeurs de la copie
sont justes à quelques unités près, avec des écarts plus grands aux rangs 14
($68709$ pour $68698$) et 22 ($54513$ pour $54525$). La seconde colonne est la
première divisée par 500, arrondie à des fractions simples ; une seule de ses
valeurs est grossièrement fausse, celle du rang 13, $143\frac34$ au lieu de
$141{,}42$, c'est-à-dire $100\sqrt2$.\footnote{La transcription note que
plusieurs de ces fractions sont petites et douteuses. La construction que la
remarque suivante attribue à l'almérie de Le Maire — onze moyennes
proportionnelles entre deux longueurs en raison double, puis une moyenne entre
chacune des onze — donne exactement cette division.}

La remarque sur la page 54 rapporte une construction de M. de Fortemaison pour
placer les touches du luth, et conclut qu'elle ne peut être juste parce que ses
douze longueurs ne sont pas « continuellement proportionnelles ». C'est le bon
critère : des touches justes pour douze demi-tons égaux placent les longueurs
vibrantes en progression géométrique de raison $2^{-1/12}$. La remarque sur la
page 68 renvoie à Descartes pour la construction des deux moyennes
proportionnelles par une parabole et un cercle ; celle sur la page 225, sur la
section en moyenne et extrême raison, rappelle justement qu'en retranchant le
petit segment du grand on retrouve la même section, et que douze sections
successives ne donnent pas douze demi-tons : le rapport de la section,
$\frac{1+\sqrt5}{2}$, n'a rien de commun avec $2^{1/12}$.

\subsection*{L'or et le laiton}

La remarque sur la page 151 soutient contre Galilée que si une corde d'or sonne
plus bas qu'une corde de laiton de même longueur, grosseur et tension, c'est
« parce qu'elle est plus molle, et non parce qu'elle est plus pesante », et celle
sur la page 152 que les deux cordes font le triton et non la quinte. Sur les deux
points, c'est Galilée qui a raison. Par la loi de la corde,
$f = \frac{1}{2L}\sqrt{T/\mu}$, des cordes de même longueur, de même section et
de même tension ont des fréquences en raison inverse de la racine de leurs
densités ; l'or ($19{,}3$) et le laiton ($8{,}4$ à $8{,}7$) donnent un rapport de
$1{,}49$ à $1{,}52$, c'est-à-dire la quinte, $1{,}5$. Le triton, $\sqrt2$,
demanderait un rapport de densités de 2.

\subsection*{Les bois de Faber, et les flûtes}

La table de Faber (page 176) est juste dans sa construction : sa deuxième
colonne, $162, 182\frac14, 192, 216, 243, 256, 288, 324$, est une gamme
pythagoricienne ; la première en donne des racines approchées ; la troisième est
une autre gamme pythagoricienne, de 9 à 18 ; la quatrième n'est pas faite des
carrés de la troisième, comme le dit le texte, mais de ces carrés divisés par 18.
La « racine imaginable » de $\frac98$ que Faber prend pour $\frac{18}{17}$ est
approchée : $17^2$ vaut 289 et non 288, et $\bigl(\frac{18}{17}\bigr)^2 =
\frac{324}{289}$ est plus petit que $\frac98$ de $\frac{1}{289}$ de sa
valeur.\footnote{Le texte écrit « 17 multiplié par soi, c'est-à-dire 288 ». La
page suivante dit que $\frac{17}{18}$ est la racine de la sesquioctave « moins
$\frac{1}{288}$ ».}

La remarque qui suit compare des flûtes de même longueur et de même tour, l'une
ronde, l'autre carrée, l'autre triangulaire. Avec un tour de 22 et
$\pi = \frac{22}{7}$, le cercle a pour aire $38\frac12$, le carré $30\frac14$,
et leurs aires sont comme 154 à 121, soit 14 à 11 : juste. Le triangle a pour
côté $7\frac13$ et pour aire, par la règle des $\frac{13}{30}$ déjà vue,
$23\frac{41}{135}$ ; l'aire vraie est $23{,}29$. Le cercle est au triangle comme
$1{,}65$ ; la copie dit « comme 77 à 46 » ($1{,}67$), ce qui revient à arrondir
l'aire du triangle à 23, et elle écrit le rapport dans le mauvais sens (« sa
capacité sera à celle du cercle »).

\subsection*{La trompette}

La remarque sur la page 250 constate que les sons de la trompette « suivent
l'ordre naturel des nombres 1, 2, 3, 4 », puis « sautent de 6 à 8 ». C'est la
série harmonique ; le septième harmonique existe, mais il est faux par rapport
à la gamme, et les trompettistes l'évitent, ce qui explique le saut.

\subsection*{La cloche}

La copie donne la construction complète du profil d'une cloche (vues 33 à 35),
avec un dessin sur un feuillet volant (vue 34), une table des longueurs et une
table des volumes. Les volumes sont cohérents : les solides des tranches moins
leurs vides donnent bien $49677$, $113443$ et $45594$ lignes cubes, les deux vides
du bas font $193176$, et le total des métaux et des filets est $233550$, comme
l'écrit la copie.

La remarque sur la page 19 calcule ensuite les angles de la construction, et la
marge note : « ces angles sont faux ». Trois le sont, par des fautes de chiffres,
les autres sont justes. Dans le triangle rectangle EOB, de côtés $OB = 35$ et
$OE = 115$ (ou d'hypoténuse $EB = 120$), l'angle en B vaut $73^\circ 4'$ et non
$75^\circ 3'$ ; l'angle en E, $16^\circ 56'$ (la copie : $16^\circ 57'$), en est
le complément, et la somme $75^\circ 3' + 16^\circ 57' = 92^\circ$ trahit la
faute. Les angles à la base du triangle isocèle RBI, $53^\circ 28'$, et l'angle
IXP, $36^\circ 32'$, sont justes, et ils sont calculés à partir de $73^\circ$ et
non de $75^\circ$. Dans le triangle isocèle PBR, de côtés 75 et de base 50,
l'angle au sommet vaut $2\arcsin\frac13 = 38^\circ 57'$ et non $38^\circ 36'$, les
angles à la base $70^\circ 32'$ et non $72^\circ 32'$ ; l'angle PBI,
$73^\circ 4' - 38^\circ 57' = 34^\circ 7'$, est juste.\footnote{La marge corrige
aussi « RBD » en « RBP », à juste titre. La table des longueurs donne $Bu =
4\frac{3}{10}$, le texte $uB = 4\frac{7}{20}$.}

Le poids du métal : quatre pouces cubes font $4 \cdot 1728 = 6912$ lignes
cubes ; ils pèsent 13053 grains, que la copie dit « 22 onces et 21 grains » — il
manque 5 gros, car $13053 = 22 \cdot 576 + 5 \cdot 72 + 21$. Une livre de ce
métal occupe donc $6912 \cdot 9216 / 13053 \approx 4880$ lignes cubes, et
$52\frac45$ livres en occupent $4880 \cdot 52\frac45 = 257664$ : juste. Le
volume d'eau déplacé dans le vase cylindrique de 13 pouces 9 lignes $\frac13$ de
diamètre ($165\frac13$ lignes), l'eau montant de 12 lignes, est
$\pi \cdot (82\frac23)^2 \cdot 12 \approx 257\,630$ lignes cubes ($257\,730$
avec $\frac{22}{7}$). Le nombre que la transcription lit, 197941, ne peut être
juste, et la marge dit en effet qu'« il y a quelque autre erreur en ce
calcul » ; l'écart avec 257664 est de quelques dizaines de lignes, non des 123
que dit la copie.

\subsection*{Les verres}

La remarque sur la page 33 conteste Kircher, selon qui un verre à moitié plein
sonnerait l'octave au-dessous du verre vide. La copie a raison de douter : la
hauteur d'un verre frotté baisse quand on le remplit, mais bien moins que d'une
octave à mi-hauteur, l'eau ne chargeant que la partie basse de la paroi.

\section*{L'utilité de l'harmonie}

\pagerange{37}{40}
\subsection*{Le poids de l'air}

Descartes a trouvé que l'air est à l'eau comme 1 à 145, avec une fiole soufflée
à la lampe qui pesait $78\frac12$ grains froide et 78 chauffée, et qui, refroidie,
aspira $72\frac12$ grains d'eau : $\frac12 : 72\frac12 = 1 : 145$. L'arithmétique
est juste, et la méthode aussi : le demi-grain perdu est le poids de l'air chassé
par la chaleur, et l'eau aspirée en mesure le volume. C'est la pesée qui ne peut
l'être : l'eau étant en réalité un peu plus de 800 fois plus dense que l'air, la
fiole n'aurait dû perdre qu'un douzième de grain environ, et une fiole chaude
se pèse mal. Les
autres valeurs du volume, 400 (Galilée), 1300, $1356\frac14$, 1380, 1700, se
répartissent autour de ce nombre ; celle de Galilée est trop faible du double.

\subsection*{Miroirs ardents et lunettes}

La remarque sur la page 28 démontre qu'on ne peut faire « de miroirs qui
brûlent jusqu'à l'infini » : les rayons venus de chaque point du soleil ne sont
pas parallèles entre eux, et le faisceau s'élargit comme le diamètre apparent
du soleil. C'est juste : l'image du soleil au foyer d'un miroir de distance
focale $f$ a un diamètre d'environ $f/107$, et la chaleur qu'on concentre
diminue quand $f$ grandit. La lunette à deux paraboles de même foyer, la petite
renvoyant les rayons de la grande « quasi parallèles », est un télescope à
réflexion afocal ; avec des ouvertures de 8 pouces et d'un demi-pouce, la
lumière reçue est multipliée par $(8 : \frac12)^2 = 256$, comme le dit la
copie.\footnote{Le nom de télescope afocal à deux paraboloïdes confocaux est
moderne.}

\subsection*{Les ellipses rationnelles}

La remarque sur la page 31 annonce un problème de nombres : le grand diamètre
d'une ellipse étant donné, trouver des ellipses dont toutes les lignes — la
perpendiculaire EF élevée au foyer, la distance AF du sommet au foyer, le petit
diamètre — soient rationnelles, avec la distance des foyers plus grande que le
petit diamètre et l'aire de l'ellipse plus grande que celle du cercle décrit
sur la distance des foyers ; et elle donne, « pour 12 ellipses et non plus », le
nombre 434788275897787529.

Si $a$, $b$, $c$ sont le demi-grand axe, le demi-petit axe et la demi-distance
des foyers, on a $a^2 = b^2 + c^2$, $AF = a - c$ et $EF = b^2/a$ : rendre ces
lignes rationnelles, c'est demander un triangle rectangle en nombres
d'hypoténuse $a$. Les deux conditions s'écrivent $c > b$ et $ab > c^2$ ; en
posant $b = a\sin\beta$, elles reviennent à $\sin\beta < \frac{\sqrt2}{2}$ et
$\sin\beta > \frac{\sqrt5 - 1}{2}$, soit $38{,}2^\circ < \beta < 45^\circ$. Le
problème est donc bien posé, et c'est un problème de représentations d'un carré
en somme de deux carrés. Mais le nombre tel que la transcription le lit se
factorise en $1051 \cdot 5987 \cdot 69098059417$, dont les deux premiers
facteurs sont de la forme $4k+3$ : il n'est hypoténuse que d'un seul triangle
rectangle en nombres entiers, qui ne satisfait d'ailleurs pas aux deux
conditions, et ne peut servir à douze ellipses. Cette lecture
ne peut dire si l'écart vient de la lecture du nombre, de la copie, ou d'un
énoncé plus fin que celui qu'elle en tire.\footnote{La transcription lit le
nombre chiffre à chiffre sur une coupe agrandie et note que le dernier chiffre,
contre le pli, est le moins sûr ; aucune des dix valeurs possibles de ce chiffre
ne donne plus de deux triangles qui satisfassent aux deux conditions. Si l'on n'exige que
des lignes rationnelles, et non entières, tout nombre convient, et une
infinité d'ellipses. La copie ajoute que « l'analyse n'a pu servir à cette
question jusqu'à présent ».}

\subsection*{Portées}

Galilée démontre que la portée à $45^\circ$ est 10000 quand celle d'un degré est
349 : c'est $1/\sin 2^\circ = 28{,}65$, juste dans le vide. La copie répond que
dans l'air la portée à $45^\circ$ n'est guère que triple de celle d'un degré, et
que la balle, qui fait plus d'effet de près que de loin, ne garde pas une vitesse
horizontale égale. Les deux sont vrais, chacun dans son milieu. La remarque sur
la page 42 porte une date : « le dernier jour de mai 1636, sur la foire à
Saint-Cloud », en tête de trois mesures de portées d'arquebuse.\footnote{La
transcription note que le premier chiffre de l'année est petit et peu formé et
marque l'année comme douteuse.} La remarque sur la page 61 réfute l'argument
d'une balance dont un poids pend au bout d'un bras, par la même raison que la
réfutation de Beaugrand lue plus loin : un poids suspendu tire suivant son fil,
et l'on ne peut composer son effet comme s'il était attaché au bras.

\section*{Les nouvelles observations physiques et mathématiques}

\pagerange{40}{44}
\subsection*{L'écoulement de l'eau}

« J'ai premièrement trouvé que les vitesses de la sortie de l'eau par un même
trou sont en raison sous-doublée des hauteurs des tuyaux » : un tuyau de 4 pieds
débite deux fois plus vite qu'un tuyau d'un pied, un tuyau de 16 pieds quatre
fois, pourvu qu'on les tienne pleins. C'est la loi qu'on appelle aujourd'hui
de Torricelli, $v = \sqrt{2gh}$.\footnote{Le nom est moderne, et cette lecture ne
dit rien de la priorité, que la copie semble revendiquer par son
« premièrement » : c'est l'affaire d'une autre enquête, sur des sources datées.}
Les mesures qui suivent se vérifient en partie. Un jet horizontal sortant à 6
pouces du sol d'un tuyau de 4 pieds doit porter, sans pertes, à
$2\sqrt{48 \cdot 6} \approx 34$ pouces : la copie mesure trois pieds, ce qui est
remarquablement juste. Le jet à $45^\circ$ devrait porter à environ 8 pieds, la
copie en mesure 6 ; le jet vertical devrait remonter à 4 pieds, la copie en
mesure un peu plus de 3, « quasi un quart » de moins : ce sont les pertes.

Le calcul de la vue 41 est juste dans son principe : pour que l'eau jaillisse
60 fois plus vite, il faut un tuyau $60^2 = 3600$ fois plus haut, soit
$12 \cdot 3600 = 43200$ pieds. La copie écrit « en raison doublée de 1 à 60, ou
de 12 à 720 », où 720 est $12 \cdot 60$ et non le terme de la raison doublée,
qui est 43200.

\subsection*{La chute dans l'eau, et une date}

« Enfin le 4\textsuperscript{e} juin 1638 j'ai trouvé que le plomb est à la vessie
comme 1890 à 1 » (vue 41) : c'est la deuxième date d'expérience du volume. La
remarque voisine rapporte « l'expérience contre Galilée » : l'eau pesant douze
fois moins que le plomb, la règle de Galilée voudrait que le plomb tombe dans
l'eau en deux secondes presque autant que dans l'air, et il ne tombe que de 12
pieds. Avec la chute de 12 pieds dans la première seconde, qui en fait 48 en
deux, la règle donne $48 \cdot \frac{11}{12} = 44$ pieds ; la copie écrit 42, et
la pièce latine de la fin, qui reprend le même argument, dit bien 12 contre
11.\footnote{La copie écrit aussi que l'air ôte au plomb la
« $\frac{1}{480}$ partie » de son mouvement ; avec l'air 400 fois plus léger que
l'eau et le plomb 12 fois plus lourd, c'est $\frac{1}{4800}$, comme le dit la
pièce latine.} La règle de Galilée est en effet fausse dans l'eau : elle ne
tient compte que de la poussée, et c'est la résistance du fluide, croissante
avec la vitesse, qui limite très vite la chute.

La remarque sur la page 10 décrit la parabole comme composée d'un mouvement
horizontal uniforme et d'une montée qui repasse, en sens inverse, par tous les
degrés de vitesse de la descente : c'est juste dans le vide. Puis la copie
ajoute : « Il faut corriger tout ceci, puisque les flèches convainquent que la
descente est plus longue à se faire que la montée. » C'est juste dans l'air.

\subsection*{Contre la géostatique de Beaugrand}

La remarque sur la page 17 réfute en cinq points une balance dont les bras
seraient dirigés vers le centre de la Terre, et d'où son auteur tirait que les
poids changent de pesanteur avec leur distance au centre. Le deuxième point est
le plus fort, et il est juste : le principe d'Archimède et de Pappus sur le
centre de gravité de deux poids, qui partage la droite de leurs centres en
raison inverse des poids, n'est vrai que si les poids tirent suivant des
lignes parallèles ; pour des forces qui concourent au centre de la Terre, il
n'y a pas en général de centre de gravité fixe, et l'auteur qui suppose le
contraire « ne prouve rien ». Le quatrième point relève une pétition de
principe exactement : pour prouver que la distance change la pesanteur,
l'auteur suppose le centre de gravité placé comme si elle ne la changeait
pas.\footnote{La copie nomme « M. de Beaugrand » et « l'auteur de la
géostatique ». Dans la gravitation newtonienne, le poids varie bien avec la
distance au centre, mais de façon insensible à l'échelle d'une balance.}

\subsection*{La neuvième observation : le tempérament égal}

La grande table de la vue 43 compare, pour la tierce mineure, la tierce majeure,
la quarte, la quinte et les deux sixtes, quatre valeurs justes ou approchées sur
une corde de 100000 parties — la « parfaite », une « approchante », une
« imparfaite », une « du tout éloignée » — à celle de « l'accord de Gallé »,
obtenue « par moyennes proportionnelles », c'est-à-dire le tempérament égal à
douze demi-tons. Voici les valeurs, recalculées :
\[
\begin{array}{l|c|c|c|c|c}
 & \text{parfaite} & \text{approchante} & \text{imparfaite}
   & \text{éloignée} & \text{tempérament égal} \\ \hline
\text{tierce mineure} & \frac56 = 83333 & \frac{1024}{1215} = 84280
   & \frac{27}{32} = 84375 & \frac{64}{75} = 85333 & 2^{-3/12} = 84090 \\
\text{tierce majeure} & \frac45 = 80000 & \frac{405}{512} = 79102
   & \frac{64}{81} = 79012 & \frac{25}{32} = 78125 & 2^{-4/12} = 79370 \\
\text{quarte} & \frac34 = 75000 & \frac{512}{675} = 75852
   & \frac{20}{27} = 74074 & & 2^{-5/12} = 74915 \\
\text{quinte} & \frac23 = 66667 & \frac{675}{1024} = 65918
   & \frac{27}{40} = 67500 & & 2^{-7/12} = 66742 \\
\text{sixte mineure} & \frac58 = 62500 & \frac{256}{405} = 63210
   & \frac{81}{128} = 63281 & \frac{16}{25} = 64000 & 2^{-8/12} = 62996 \\
\text{sixte majeure} & \frac35 = 60000 & \frac{1215}{2048} = 59326
   & \frac{16}{27} = 59259 & \frac{75}{128} = 58594 & 2^{-9/12} = 59460
\end{array}
\]
Toutes les valeurs de la copie sont justes, à trois exceptions près :
$\frac{75}{128}$ vaut 58594 et non 58994 ; $\frac{81}{128}$ vaut $63281\frac14$,
écrit 63282 ; et la dernière ligne, intitulée « Sexta minor » comme la
précédente, est celle de la sixte majeure. Les « différences avec la parfaite »
ne sont pas des différences de longueurs mais des écarts relatifs, en cent
millièmes : 1123 est l'écart de $\frac{2048}{2025}$, 1235 celui du comma
syntonique $\frac{81}{80}$, 2344 celui du dièse $\frac{128}{125}$ ; et pour le
tempérament égal, 113 pour les quintes et les quartes, 789 pour les tierces
majeures et les sixtes mineures, 900 pour les tierces mineures et les sixtes
majeures. Tous sont justes.

Le paragraphe qui suit tire de là que, dans le tempérament égal, la tierce
majeure est plus près de la juste que la mineure, et il le vérifie en commas :
l'octave valant 53 commas, le ton égal vaut $8\frac56$, le demi-ton
$4\frac{5}{12}$, la tierce mineure égale $13\frac14$ au lieu de 14, trop basse
de $\frac34$ de comma, la tierce majeure $17\frac23$ au lieu de 17, trop haute de
$\frac23$ : la mineure s'écarte davantage, de $\frac{1}{12}$ de comma. C'est
exact, et conforme à la table ($900 > 789$) : en cents, $-15{,}6$ contre
$+13{,}7$. La copie en conclut que les facteurs d'orgues, qui jugent le
contraire, ont l'oreille « pervertie par accoutumance ».\footnote{« Gallé »,
avec un trait sur le e, est le nom tel que la transcription le lit ; il abrège
aussi l'en-tête de la dernière colonne, « Accord de Ga\dots ». Cette lecture ne
l'identifie pas. La fraction qui suit « chaque demi-ton 4 » et celle qui suit
« abaissée de » sont cachées dans le pli ; le calcul les donne, $\frac{5}{12}$ et
$\frac34$.}

\subsection*{La roulette}

La remarque sur la page 25 annonce qu'on a trouvé l'aire de la roulette pour
tout rapport entre la circonférence du cercle roulant et le chemin qu'il fait,
ses tangentes, ses centres de gravité, et que le solide qu'elle engendre en
tournant sur sa base est au cylindre de même hauteur « comme 8 à 5 ». Le
rapport est juste, lu dans l'autre sens : le solide engendré par une arche de
cycloïde tournant autour de sa base vaut $5\pi^2 r^3$, le cylindre circonscrit
$8\pi^2 r^3$ ; le cylindre est au solide comme 8 à 5.\footnote{« Roulette » :
la cycloïde. La copie renvoie aux lettres de Descartes pour les démonstrations,
et ne donne pas les exemples qu'elle annonce (« dont je mets ici quelques
exemples »).}

\subsection*{Les nombres amiables}

La remarque sur la page 26, en latin, donne une règle « ex Cartesio » : si
$x$ est une puissance de 2 telle que $3x - 1$, $6x - 1$ et $18x^2 - 1$ soient
premiers, les nombres $2x(18x^2 - 1)$ et $2x(3x-1)(6x-1)$ sont amiables — chacun
est la somme des parties aliquotes de l'autre. La règle est exacte ; c'est
celle qu'on nomme aujourd'hui règle de Thābit ibn Qurra. Pour $x = 2$, 8 et 64
elle donne, avec les premiers $5, 11, 71$ ; $23, 47, 1151$ ; $191, 383, 73727$,
les trois paires de la copie, $220$ et $284$, $17296$ et $18416$, $9363584$ et
$9437056$, et cette lecture a vérifié que chacune est amiable.\footnote{Si
$s(n)$ désigne la somme des diviseurs propres de $n$, on a par exemple
$s(9437056) = 9363584$ et $s(9363584) = 9437056$. La règle est ici attribuée à
Descartes par la copie ; le nom de Thābit est un usage moderne, et la question de
savoir qui l'a trouvée ou retrouvée n'est pas l'affaire de cette lecture.} Mais la
phrase finale, « et l'on en peut trouver une infinité d'autres de la même
manière », n'est pas démontrée, et elle est fausse autant qu'on le sache : on ne
connaît aucune autre valeur de $x$ pour laquelle les trois nombres soient
premiers.

\subsection*{Les nombres multiparfaits}

La treizième observation dresse des listes. Dans la langue de la copie, un nombre
dont les parties aliquotes font « le double » est un nombre dont la somme des
diviseurs propres est double de lui, c'est-à-dire, en notant $\sigma(n)$ la
somme de tous ses diviseurs, $\sigma(n) = 3n$. Les six nombres donnés —
$120$, $672$, $523776$, $459818240$, $1476304896$, $51001180160$ — sont bien tels, et ce
sont exactement les six nombres de cette sorte connus aujourd'hui. Que ce soient
les seuls, comme l'affirme la copie (« il n'y en a dans tous les nombres que les
6 qui suivent »), est une question ouverte.\footnote{En langage moderne, ce sont
les six nombres 3-parfaits connus. Aucun autre n'a été trouvé, et l'on ne sait
pas démontrer qu'il n'en existe pas.}

Les neuf nombres dont les parties aliquotes font « le triple » ($\sigma(n) =
4n$) — $30240$, $32760$, $23569920$, $45532800$, $142990848$, $1379454720$, $43861478400$, $66433720320$, $403031236608$ — le sont tous ; mais la liste n'est pas complète
dans son intervalle, car $2178540$ et $153003540480$ le sont aussi. La règle
qui passe de $459818240$ à $1379454720$ en multipliant par 3 est juste, et c'est
le cas particulier d'un lemme que le cadre de la même page énonce en général :
si $\sigma(n) = kn$ avec $k$ premier ne divisant pas $n$, alors
$\sigma(kn) = (k+1)\,kn$, puisque $\sigma$ est multiplicative et
$\sigma(k) = k + 1$. Ainsi un « sextuple » non divisible par 7, multiplié par 7,
donne un « septuple ».\footnote{Le chiffre douteux du dernier des neuf nombres,
lu « 0 ou 8 », est 0 : $403031236608$ est 4-parfait, $483031236608$ ne l'est
pas.}

Des nombres dont les parties font « le quadruple » ($\sigma(n) = 5n$), la copie
donne $14182439040$ et $30823866178560$, qui le sont, et $508666803200$, qui ne
l'est pas tel que la transcription le lit ; $518666803200$, qui n'en diffère que
par un chiffre, l'est.\footnote{Ce qui fait penser à une faute de lecture ou de
copie d'un seul chiffre ; la page doit trancher.} Le cadre donne enfin un produit de
vingt et un facteurs, $2^{31}$, $13$, $17$, $31$, $41$, $43$, $61$, $83$, $163$,
$257$, $307$, $331$, $467$, $613$, $1093$, $11^3$, $7^4$, $2801$, $65537$, $547^2$
et « 594323 », qui ferait un « quadruple » ; tel qu'il est lu, il n'en est pas un
($\sigma/n \approx 3{,}33$). Mais si l'on remplace 594323 par
$1594323 = 3^{13}$, le produit $P$ vérifie exactement $\sigma(P) = 5P$, et,
n'étant pas divisible par 5, il donne $\sigma(5P) = 6 \cdot 5P$, le
« quintuple » que le cadre annonce. Le premier chiffre du dernier facteur
manque, sur la page ou dans la lecture.\footnote{Le nombre qu'il faudrait est
déterminé : le rapport $\sigma(x)/x$ que doit avoir le dernier facteur pour que le
produit soit 5-parfait vaut $\frac{2391484}{1594323}$, et $2391484 =
\sigma(3^{13})$. Le calcul est de cette lecture.}

\section*{Les feuillets blancs de la fin}

\pagerange{45}{52}
Sous la rubrique « Au papier blanc après tous les livres de l'harmonie »,
huit pages de nombres, sans renvoi au livre sauf une fois.

\subsection*{Les carrés magiques emboîtés}

Un carré magique d'ordre $n$ est une disposition des nombres $1$ à $n^2$ en carré
telle que chaque rangée, chaque colonne et chaque diagonale ait la même somme,
qui est alors $\frac{n(n^2+1)}{2}$. Le carré de 22 de la vue 45 a pour somme
$22 \cdot 485 / 2 = 5335$. Il contient, emboîtés au centre, des carrés de 16, 12 et
10 qui sont eux-mêmes magiques : ce sont des carrés « à enceintes », qu'on peut
dépouiller de leurs bordures en gardant la propriété. Tel que la transcription
le lit, toutes ses colonnes et ses deux diagonales font 5335 ; les carrés de 12
(rangées et colonnes 6 à 17, de coins 241, 242, 243, 244) et de 10 (rangées 7 à
16, de coins 33, 46, 452, 439) sont exactement magiques, de sommes 2910 et 2425 ;
seules les rangées 18 et 19 font 5338 et 5332. La copie a tracé deux traits entre
les cases 433 et 430 de ces rangées ; les échanger corrige les rangées, mais
dérange alors les colonnes 18 et 19 de trois unités. Un autre chiffre du bloc est
donc probablement mal lu ou mal copié.\footnote{La transcription écrit que
l'échange « rétablit les deux sommes », ce qui n'est vrai que des rangées. Trois
cases tachées y sont restituées par la somme : 254, 256 et 154, ce dernier écrit
en clair dans la marge, en face de la rangée 21.}

Les notes de marge sont les plus intéressantes. « Chaque rang de 16 contient
2056 selon le Père Mersenne, mais il est faux, car il contient 3880 » ; de même
870 et 2910 pour le carré de 12, 505 et 2425 pour celui de 10. Les deux nombres
de chaque paire sont justes, de deux objets différents. $2056 =
\frac{16 \cdot 257}{2}$, $870 = \frac{12 \cdot 145}{2}$ et $505 =
\frac{10 \cdot 101}{2}$ sont les sommes magiques de carrés d'ordre 16, 12 et 10
faits des nombres $1$ à $n^2$ ; $3880$, $2910$ et $2425$ sont celles des carrés
intérieurs du carré de 22, qui sont faits de nombres pris entre 1 et 484, et dont
la somme est $k \cdot \frac{485}{2}$. La correction de la copie est juste pour le
carré qu'elle a sous les yeux.

Les carrés de 9 et de 8 « pour les planètes » sont magiques (sommes 369 et 260), et
à enceintes eux aussi : le carré de 9 contient des carrés magiques de 7, 5 et 3
(sommes 287, 205, 123), attribués à Vénus, Mars et Saturne, le carré de 8 des
carrés de 6 et 4 (sommes 195 et 130), pour le Soleil et Jupiter. Les deux carrés
de 4 des vues 45 et 46 sont magiques. Que le carré de 4 se varie « en 880
façons » sans compter les huit retournements, et que le carré de 3 ne se fasse
que d'une, est exact : il y a 880 carrés magiques d'ordre 4 à rotation et
symétrie près, et un seul d'ordre 3. Que « toutes ses dispositions différentes »
soient $16! = 20922789888000$ est juste aussi. La vue 46 ajoute que, parmi les 880,
« il y en a 16 seulement » qui ont aux quatre coins quatre nombres consécutifs, et
que chacune se voit en 8 manières, soit 128 : cette lecture a compté les carrés
magiques d'ordre 4 et trouve bien 128 carrés, 16 à la symétrie près. Mais
l'exemple que donne la copie, « comme est 4, 5, 6, 7 », est impossible : les
quatre coins d'un carré magique d'ordre 4 font toujours 34, et quatre nombres
consécutifs de somme 34 ne peuvent être que 7, 8, 9, 10.\footnote{Le carré de 14
« d'une lettre de M. Fermat au Père Mersenne du premier [?] avril 1640 », dont
la copie donne les sommes 1379, 985 et 591 pour les carrés de 14, 10 et 6, n'est
pas recopié ; les trois sommes sont justes, $k \cdot \frac{197}{2}$. La
transcription lit la date par inférence, le jour étant un signe souligné non
lu. Le nombre $6742112993280000$ des variations du carré de 9 vaut
$2^{25} \cdot 3^8 \cdot 5^4 \cdot 7^2 = 8 \cdot (8! \cdot 6!)^2$ ; cette lecture
ne sait pas le rattacher à la description qu'en donne la copie.}

\subsection*{Deux nombres dont on demande s'ils sont premiers}

« On propose aux arithméticiens le nombre 2027651281 : qu'ils disent en quatre
jours s'il est premier ou composé ; s'ils ne le peuvent, nous le donnerons en une
heure par des règles certaines et générales. » Il est composé :
$2027651281 = 44021 \cdot 46061$. La factorisation se trouve en une heure, en
effet, par la méthode qui porte aujourd'hui le nom de Fermat : on cherche un
carré $x^2$ tel que $x^2 - N$ soit un carré, en partant de $\lceil\sqrt N\,\rceil
= 45030$ ; au douzième essai, $45041^2 - 2027651281 = 1020^2$, et
$N = (45041 - 1020)(45041 + 1020)$.\footnote{Le nom de méthode de Fermat est
moderne. La copie ne dit pas de qui est le défi.}

Le second est en latin, sur la vue 50 : « on ne trouvera pas par la voie
commune si 1073602561 est premier ou non, à moins de le diviser par tous les
premiers jusqu'à 8191 ». Il ne l'est pas, et 8191 en est justement le plus petit
diviseur : $1073602561 = 8191 \cdot 131071 = (2^{13} - 1)(2^{17} - 1)$, produit de
deux nombres premiers de Mersenne. Diviser jusqu'à 8191 trouve donc le facteur ;
mais pour prouver qu'un nombre de cette taille est premier, il faudrait
diviser jusqu'à sa racine, 32766.\footnote{« Viam communem non reperies » est la
lecture développée par la transcription de « viã coĩ nõ reperiz ». Le nom de
nombres de Mersenne, pour les $2^p - 1$, est moderne.}

La même page juge d'autres nombres : « 1373 est premier, et non 1369, carré de
37, ni 1371 » — juste, $1371 = 3 \cdot 457$. Mais « 6679 n'est pas premier, il se
divise par 23 et 29 » est faux : $23 \cdot 29 = 667$, et 6679 est premier.

\subsection*{Compter les parties aliquotes}

Les « parties aliquotes » d'un nombre sont ses diviseurs autres que lui-même ;
leur nombre est $\tau(n) - 1$, où $\tau(n)$ compte tous les diviseurs, et si
$n = p_1^{e_1}\cdots p_r^{e_r}$, alors $\tau(n) = (e_1+1)\cdots(e_r+1)$. C'est
exactement la première règle de la vue 46 : pour 108, produit du carré 4 et du
cube 27, on prend les exposants 2 et 3, on ajoute un à chacun, $3 \cdot 4 = 12$,
et l'on ôte un : 11 parties aliquotes. La « seconde méthode » —
$2 \cdot 3 + 2 + 3 = 11$ — est l'identité $(a+1)(b+1) - 1 = ab + a + b$ et ne vaut
que pour deux facteurs. Les exemples de la copie sont justes : 20 et 12 et 45
en ont 5, 360 en a 23, 400 en a 14, 5000 en a 19, 12000 en a 47, 888 en a 15,
666 en a 11, 6000 en a 39, et 5040 en a 59.

Deux sont faux. 15000 n'a pas « 80 parties moins une » mais 39 :
$15000 = 2^3 \cdot 3 \cdot 5^4$, et $\tau = 4 \cdot 2 \cdot 5 = 40$ ; la copie
donne au facteur 3 un exposant 3. Et la règle inverse — trouver le moindre
nombre qui a un nombre donné de parties aliquotes — est fausse sur l'exemple
même de la copie. Pour 11 parties, il faut $\tau(n) = 12$ ; la copie décompose
$12 = 3 \cdot 4$, prend un carré et un cube, et trouve $9 \cdot 8 = 72$. Mais
12 se décompose aussi en $2 \cdot 2 \cdot 3$, ce qui donne
$2^2 \cdot 3 \cdot 5 = 60$, plus petit, et 60 a bien 11 parties
aliquotes.\footnote{Les formes possibles pour $\tau(n) = 12$ sont $p^{11}$,
$p^5 q$, $p^3 q^2$ et $p^2 q r$, dont les plus petites valeurs sont 2048, 96, 72
et 60. Pour 14 parties, le moindre est bien $16 \cdot 9 = 144$, comme le dit la
copie ; pour 59, c'est bien 5040, le « nombre civil de Platon », comme le dit la
page latine (« nec ullus minor eas habet »).}

\subsection*{Sommer les parties aliquotes}

La quatrième règle (vues 47 et 48) calcule la somme des parties aliquotes par la
multiplicativité de $\sigma$ : on décompose le nombre en puissances de premiers,
on prend pour chacune la somme de ses diviseurs, on multiplie, et l'on ôte le
nombre. La règle est juste et la plupart des exemples aussi : 12 donne 16, 20
donne 22, 5040 donne 14304 ($31 \cdot 13 \cdot 8 \cdot 6 = 19344$), 12000 donne
27312, 5000 donne 6715, 888 donne 1392, 6000 donne 13344. Deux sont faux. Pour
666, le produit $3 \cdot 38 \cdot 13$ vaut 1482 et non 1282, et la somme des
parties est 816 et non 697. Pour 560400, la somme est 1238592, et non 1374480 —
ni pour 500400, l'autre lecture du chiffre douteux.\footnote{On a $560400 =
2^4 \cdot 3 \cdot 5^2 \cdot 467$. La table des trente-neuf parties aliquotes de
6000 (vue 50) est juste si l'on lit 50 et 250 aux deux places où la transcription
lit « 5 » et « 120 », nombres qui sont déjà ailleurs dans la table.}

La troisième règle est l'inverse : trouver un nombre dont la somme des parties
est donnée. Si $s - 1$ est un nombre premier $p$, le carré $p^2$ a pour parties
$1$ et $p$ et convient ; si $s - 1 = p + q$ avec $p$ et $q$ premiers distincts,
$pq$ convient. Les exemples sont justes : 121 pour 12, qui est bien le seul ;
529 pour 24, le seul aussi ; 95 pour 25. Mais 95 n'est pas le seul pour 25 : 119
et 143 conviennent aussi, puisque $24 = 5 + 19 = 7 + 17 = 11 + 13$.

\subsection*{Les triplets en proportion continue}

« Le moindre nombre de ceux qui contiennent vingt-quatre fois, et non plus, trois
nombres proportionnels, est 24843, et non 753571, qui les contient aussi. » Le
sens est donné par l'exemple : 91 « contient » les quatre triplets $1, 9, 81$ ;
$25, 30, 36$ ; $13, 26, 52$ ; $7, 21, 63$, qui sont en progression géométrique et
ont chacun pour somme 91. Écrire $N = a + b + c$ avec $b^2 = ac$, c'est écrire
$N = d(x^2 + xy + y^2)$ avec $a = dy^2$, $b = dxy$, $c = dx^2$ et $x, y$ premiers
entre eux, et les nombres premiers qui entrent dans $x^2 + xy + y^2$ sont 3 et
ceux de la forme $6k+1$ — d'où la condition de la copie que les premiers
« doivent être 6 plus un ». Cette lecture a compté ces triplets pour tous les
nombres jusqu'à 800000, en écartant le triplet trivial $a = b = c$. Les
affirmations de la copie sont exactes : 91 en a 4 et c'est le moindre ;
$24843 = 3 \cdot 7^2 \cdot 13^2$ en a 24 et c'est le moindre ; $753571 =
7^3 \cdot 13^3$ en a 24 aussi. Seule la note de marge se trompe, qui donne 49
pour le moindre nombre à deux triplets : 21 en a deux, $1 + 4 + 16$ et
$3 + 6 + 12$.\footnote{Pour $N$ premier à 3, le nombre de triplets est
$\frac12\bigl(\prod(2e_i + 1) - 1\bigr)$, le produit portant sur les exposants
des premiers de la forme $6k+1$ ; un facteur 3 le double. D'où 4 pour
$7 \cdot 13$, 24 pour $7^3 \cdot 13^3$, et $2 \cdot 12$ pour
$3 \cdot 7^2 \cdot 13^2$.}

\subsection*{Critères de divisibilité}

Un nombre qui finit par 0 ou 5 est divisible par 5 ; un nombre dont la somme des
chiffres, les 9 ôtés, laisse 3, 6 ou 9 est divisible par 3 ; s'il ne reste rien,
par 9. Exemple juste : les chiffres de $75512713527351$ font 54, et le nombre est
divisible par 9.

\subsection*{Sommes de puissances, en latin}

La page latine de la vue 50 et la suivante donnent des règles justes. Le
triangulaire d'un nombre impair est ce nombre multiplié par son milieu
($T_7 = 7 \cdot 4 = 28$), d'un nombre pair sa moitié multipliée par le suivant
($T_6 = 3 \cdot 7 = 21$). La somme des $n$ premiers triangulaires est $n$
multiplié par le tiers du triangulaire suivant : $\sum_{k \le n} T_k =
\frac{n\,T_{n+1}}{3} = \frac{n(n+1)(n+2)}{6}$, soit 35 pour $n = 5$. La somme
des $n$ premiers cubes est le carré du $n$-ième triangulaire : $465^2 = 216225$
pour $n = 30$, et $10^2 = 100$ pour $n = 4$. La somme des $n$ premiers carrés est
le double de $n$ fois le tiers du triangulaire suivant, moins le triangulaire de
$n$ : $2 \cdot \frac{n\,T_{n+1}}{3} - T_n = \frac{n(n+1)(2n+1)}{6}$, soit 30 pour
$n = 4$ et 9455 pour $n = 30$. Tout est exact.

Les quatre premiers « problèmes des pierres » — des pierres posées le long d'un
chemin, qu'on rapporte une à une au point de départ — sont justes aussi. Vingt
pierres à un pas l'une de l'autre demandent $20 \cdot 21 = 420$ pas ; cent
pierres aux intervalles $1, 2, 3, \dots$ demandent $2 \cdot \frac{100 \cdot 101
\cdot 102}{6} = 343400$ pas ; douze pierres aux intervalles $1, 2, 4, \dots$
sont aux distances $2^k - 1$, dont la somme est $2^{13} - 2 - 12 = 8178$, soit
16356 pas ; cinq pierres aux intervalles $1$, $3$, $9$, $27$, $81$ sont aux distances
$1$, $4$, $13$, $40$, $121$, de somme 179, soit 358 pas. Le cinquième problème est trop
incertain dans sa lettre pour être vérifié.\footnote{La transcription le dit :
« non jubet se sedere » et « sint in 4\textsuperscript{o} praecedenti puncto »
sont lus tels qu'ils sont écrits, sans que le sens soit établi.}

\subsection*{Différences de puissances}

La différence de deux cubes consécutifs est un nombre hexagonal centré :
$(n+1)^3 - n^3 = 3n^2 + 3n + 1 = 6T_n + 1$ ; ainsi $64 - 27 = 6 \cdot 6 + 1 = 37$.
Pour les quatrièmes puissances, la copie ajoute au cube du moindre nombre la
« colonne » de l'hexagone central du plus grand : $4^4 - 3^4 = 4 \cdot 37 + 27
= 175$, ce qui est l'identité $n^4 - (n-1)^4 = n\bigl(n^3 - (n-1)^3\bigr) +
(n-1)^3$. Pour les cinquièmes, la règle est plus curieuse et exacte : on double le
triangulaire du moindre nombre $m$, on prend le triangulaire de ce double, et
l'on écrit un 1 à la suite. Pour $m = 3$, $T_3 = 6$, $T_{12} = 78$, et l'on a
781, qui est $4^5 - 3^5$. En général $10\,T_{m(m+1)} + 1 = 5m^4 + 10m^3 + 10m^2 +
5m + 1 = (m+1)^5 - m^5$.\footnote{« Quarré quarré » : la quatrième puissance ;
« Q. C. » : la cinquième, quarré-cube. La liste des premières quatrièmes
puissances saute $6^4 = 1296$, et sa dernière différence, 1776, est
$7^4 - 5^4$.}

\subsection*{Les dés}

Trois dés donnent $6^3 = 216$ coups, et les points de 3 à 10 s'obtiennent de
$1$, $3$, $6$, $10$, $15$, $21$, $25$, $27$ façons — symétriquement de 18 à 11 — dont la somme
108 est la moitié de 216 : tout est exact. Les six premiers nombres sont bien
les triangulaires ; les deux suivants, 25 et 27, ne le sont plus, et la copie
les donne sans règle.

Pour deux dés, la copie se trompe. Le texte dit « 62 façons » et « 31 jeux
différents » ; il y en a 36 et 21. Le petit feuillet volant de la vue 49 montre
d'où vient la faute : le copiste a d'abord écrit, pour chaque point de 2 à 12,
toutes les décompositions en deux nombres, y compris celles où entre un nombre
plus grand que six, puis il a biffé celles-ci ; et il compte deux fois chaque
coup double (« 1 1, 1 1 »). La table de la marge (vue 51), faite après la
correction, trouve ainsi 42 coups, soit $36 + 6$, et 21 jeux sans ordre, ce
dernier nombre juste.

\section*{L'Académie parisienne aux amis de Galilée}

\pagerange{53}{55}
La dernière pièce du volume est d'un autre genre : un texte latin suivi,
intitulé « L'Académie parisienne prie les très illustres hommes, familiers et
amis de Galilée, Lyncéens, de répondre aux notes suivantes sur les livres des
dialogues ». Quinze objections numérotées renvoient aux pages des dialogues, de
la page 2 à la page 288 ; la fin explique que « ce sont les choses que les
Parisiens avaient écrites à Galilée lui-même quand il vivait encore », que sa mort
a prévenu l'envoi, et prie ses amis — Toscans, Romains, Bolonais, Génois — d'envoyer
leurs réponses à « l'éminent géomètre génois D. Santini », qui les fera
parvenir. La pièce est datée : « Lutetiae Calendis Julii anni 1643 », à Paris,
le 1\textsuperscript{er} juillet 1643. Elle ne porte pas de nom.\footnote{Le
contenu des objections — la résistance des machines selon leur grandeur, le
vide, la roue d'Aristote, les miroirs d'Archimède, le fil d'or étiré, les cordes
d'or, le prisme comme levier, l'isochronisme des oscillations, les plans
inclinés, les projectiles — est celui des \emph{Discorsi} de Galilée (1638) ;
l'identification est de cette lecture, la pièce disant seulement « dialogorum
libri ». Galilée était mort en janvier 1642.}

Les objections sont de physique, et le temps a donné raison tantôt à l'Académie,
tantôt à Galilée. Elle a raison contre lui quand elle attribue à la pression de
l'air, et non à la force du vide, l'adhérence de deux plaques polies et la limite
des pompes aspirantes (objection 2) ; quand elle nie que les petits cercles de la
roue d'Aristote laissent des « petits espaces vides » et dit qu'ils vont
seulement plus lentement (objection 5) — le petit cercle roule et glisse à la
fois ; quand elle doute que toutes les oscillations d'un même pendule soient de
durée égale (objection 11) — elles ne le sont que pour les petites amplitudes ;
et quand elle oppose l'expérience à la règle qui fait retrancher à la vitesse
d'un corps la fraction de son poids que lui ôte le milieu (objection 15), la
règle ignorant la résistance du fluide. Elle a tort quand elle nie que les corps
pèsent dans le vide (objection 8) ; quand elle attribue la voix plus grave des
cordes d'or à leur mollesse et non à leur poids (objection 9), comme on l'a vu
plus haut ; et quand elle tire de l'expérience de l'arquebuse une objection
contre la composition des mouvements (objection 14) : une balle tirée
horizontalement tombe bel et bien de quelque 15 pieds en une seconde, et si
elle touche le but, c'est que la ligne de mire et l'axe du canon ne sont pas
parallèles. Le rapport de 28 à 1 entre les portées à $45^\circ$ et à $1^\circ$,
qu'elle récuse au nom de l'expérience, est juste dans le vide et faux dans
l'air, comme l'avait déjà dit la copie.\footnote{L'objection 4 tient pour
impossibles les miroirs ardents d'Archimède. Buffon, en 1747, enflamma du bois à
plusieurs dizaines de mètres avec un assemblage de miroirs plans ; la question
de savoir ce qu'Archimède fit à Syracuse reste une autre affaire. Dans
l'objection 15, les nombres sont justes : l'air 400 fois plus léger que l'eau,
le plomb 12 fois plus lourd, donc 4800 fois plus lourd que l'air, et 12 pieds
contre 11 selon la règle.}

La pièce finit par deux demandes précises : le poids et la grandeur des
fioles avec lesquelles Galilée avait pesé l'air, « afin que nous imitions au plus
près ses observations », et le traité de la force de percussion auquel il avait
tant travaillé.

Sous la date, un tableau occupe la moitié de la page : les sept modes des Grecs,
de l'hypodorien au mixolydien, dessinés comme sept échelles en escalier, avec une
légende en français : si l'on commence les six autres modes à la mèse de
l'hypodorien, on entend « les sept tons et les sept modes des anciens, tous
diversifiés et chacun en son lieu propre, selon M. Doni ».

\section*{Ce que ces feuillets ne disent pas}

\pagerange{1}{56}
Il faut finir par les limites, parce que le titre du volume et la nature de
la copie invitent à en dire plus que les feuillets.

\textbf{Ce n'est pas l'exemplaire annoté de Mersenne.} C'est une copie de ses
notes, dans une main qui n'est pas la sienne et qui ne se nomme pas. Cette
lecture ne sait rien de l'exemplaire copié, ni s'il existe encore.

\textbf{L'attribution est celle du titre.} Les remarques sont de Mersenne parce
que la copie le dit ; trois passages le nomment à la troisième personne, dont un
qui le corrige, et cette lecture ne sait pas qui les a écrits. Les trois
feuillets de musique de devant sont d'une autre main, ne renvoient pas au
livre, et rien ne les rattache aux remarques. La pièce latine de 1643 est de la
main de la copie ; elle parle au nom de l'« Académie parisienne » et n'a pas
d'auteur nommé.

\textbf{Les renvois au livre n'ont pas été vérifiés.} Toutes les pages de
l'\emph{Harmonie} sont citées comme la copie les donne ; c'est pourquoi deux des
corrections discutées ici — celle de la page 139, et le nombre de la page 110 —
restent suspendues à une lecture du livre.

\textbf{Aucune date n'est donnée à la copie.} Les seules dates sont celles que les
feuillets écrivent : 1563 pour la « triple musique » de Le Fèvre, 1484 pour
Gafurius d'après Faber, le dernier jour de mai 1636 (l'année douteuse), 1638
pour l'observation de M. de Villiers, le 4 juin 1638, avril 1640 pour la lettre
de Fermat, le 1\textsuperscript{er} juillet 1643. Plusieurs remarques citent des
livres que la copie nomme — « l'hydraulique pneumatique », « les réflexions
physicomathématiques », « les nouvelles pensées de Galilée » que « j'ai mis en
français » — et qui répondent, par leurs titres, à des ouvrages de Mersenne
postérieurs à l'\emph{Harmonie} ; si c'est bien d'eux qu'il s'agit, les notes ont
été écrites pendant plusieurs années après la parution du livre. C'est une
inférence de cette lecture, non une date.

\textbf{Un manuscrit est cité et n'est pas ici.} « Le gros volume manuscrit
entier de cette matière », où se trouve la combinaison des huit notes de
l'octave (vues 20 et 24), n'est pas nommé autrement et n'est pas dans ce
volume.

\end{document}
